\documentclass[11pt,a4paper]{article}
\usepackage[T1]{fontenc}
\usepackage[utf8]{inputenc}
\usepackage{lmodern}
\usepackage{amsmath,amssymb,amsthm,mathrsfs}
\usepackage[margin=28mm]{geometry}
\usepackage{microtype}
\usepackage{xcolor}
\usepackage{booktabs,longtable}
\usepackage{hyperref}
\hypersetup{colorlinks=true,linkcolor=blue!45!black,urlcolor=blue!45!black,
 citecolor=blue!45!black,pdftitle={Chow Obstructions and Prescribed Actual Green Pairs}}
\newtheorem{theorem}{Theorem}[section]
\newtheorem{lemma}[theorem]{Lemma}
\newtheorem{proposition}[theorem]{Proposition}
\newtheorem{corollary}[theorem]{Corollary}
\theoremstyle{definition}
\newtheorem{definition}[theorem]{Definition}
\theoremstyle{remark}
\newtheorem{remark}[theorem]{Remark}
\newcommand{\R}{\mathbf R}
\newcommand{\B}{\mathcal B}
\newcommand{\an}{\mathrm{an}}
\newcommand{\Tr}{\operatorname{Tr}}
\newcommand{\Supp}{\operatorname{Supp}}
\newcommand{\dist}{\operatorname{dist}}
\newcommand{\Id}{\operatorname{id}}
\setlength{\emergencystretch}{2em}
\title{Chow Obstructions and Prescribed Actual Green Pairs}
\author{}
\date{23 November 2026}
\begin{document}
\maketitle
\begin{abstract}
We construct a rational-linear Chow obstruction with values in the actual
Bott--Chern quotient, and an actual curved Green pair for every prescribed
rational cycle on a smooth integral projective variety over a finite
extension of $\mathbf Q_p$. The proof retains the original valued field,
the exact prescribed support complement, and all compact actual tests.
Its inputs are actual model Chern forms, a finite trace on the
scalar--curvature module, and supported perfect complexes. We give explicit
arguments for the actual/Bedford--Taylor comparison, detection by model
Hessians, scalar PL descent with a uniform parameter, and the coordinate
Gauss product. The further vanishing of the obstruction on homologically
trivial Chow classes is a separate problem.
\end{abstract}

\paragraph*{Relation to the Strong Green programme.}
Theorem~\ref{h22:main} proves the Chow-obstruction and prescribed
curved-Green-pair construction formerly stated as Hypothesis 2.2 in
\href{https://openmath-atlas.org/projects/strong-green/papers/strong-g-programme/\#hyp:green-pairs}
{the Strong Green programme}. Its scope is the original finite extension
$K/\mathbf Q_p$, every smooth integral projective $K$-variety, the exact
prescribed support, and the full space of compact actual tests.
The argument below supplies the model-Chern construction, the minimal
finite trace, the supported-complex formula, the global primitive, and
the integral-scope extension. General vanishing on homologically trivial
Chow classes remains a separate assertion.

\tableofcontents
\newpage
\section{Chow classes and prescribed curved Green pairs}
\label{h22:main-section}

Fix a finite extension $K/\mathbf Q_p$ with its given absolute value,
and write $R=K^\circ$. For a smooth integral projective $K$-variety $X$
of dimension $n$, let $\mathcal B_X$ denote the ordinary sheaf of actual
delta-forms and put $D=d'd''$. Set
\[
 H^{r,r}_{\mathrm{BC,act}}(X)=
 \frac{\{\omega\in\mathcal B^{r,r}(X^{\mathrm{an}}):
                  d'\omega=d''\omega=0\}}
      {D\mathcal B^{r-1,r-1}(X^{\mathrm{an}})}.
\]
This is the algebraic quotient of the indicated spaces of global forms.
For a rational cycle $Z$, its support is the union of its distinct
integral components with nonzero coefficient, and
$U_Z=X^{\mathrm{an}}\setminus |Z|^{\mathrm{an}}$.

\begin{theorem}[Chow obstruction and prescribed Green pairs]
\label{h22:main}
For every such $X$ and $1\le r\le n$, there is a rational-linear map
\[
 \mathfrak o^r_{X,K}:\operatorname{CH}^r(X)_{\mathbf Q}
                 \longrightarrow H^{r,r}_{\mathrm{BC,act}}(X)
\]
with the following properties.
\begin{enumerate}
\item If $\xi$ is a rational virtual vector bundle on an integral
projective flat $R$-model of $X$ and
$\operatorname{ch}_r(\xi|_X)=a$, then
\[
              \mathfrak o^r_{X,K}(a)
                         =[\operatorname{ch}_r(\bar\xi)].
\]
The right side is independent of the model, the model extensions and
the virtual-bundle presentation. For a $K$-morphism $f:X\to Y$
between smooth integral projective varieties,
\[
 (f^{\mathrm{an}})^*\mathfrak o^r_{Y,K}(a)
                    =\mathfrak o^r_{X,K}(f^*a).
\]
The groups in degrees beyond the dimension are interpreted as zero.
\item For every prescribed rational codimension-$r$ cycle $Z$ on $X$
there are actual forms
\[
 G_Z\in\mathcal B^{r-1,r-1}(U_Z),\qquad
 \Omega_Z\in\mathcal B^{r,r}(X^{\mathrm{an}})
\]
with $d'\Omega_Z=d''\Omega_Z=0$ and
$[\Omega_Z]=\mathfrak o^r_{X,K}([Z])$, such that
\begin{equation}
\label{h22:curved-identity}
 \int_{U_Z}G_Z\wedge D\eta+\int_{Z^{\mathrm{an}}}\eta
       =\int_{X^{\mathrm{an}}}\Omega_Z\wedge\eta,
 \qquad \eta\in\mathcal B_c^{n-r,n-r}(X^{\mathrm{an}}).
\end{equation}
All products, pullbacks and integrals in this identity are in the
actual calculus over the original valued field $K$.
\end{enumerate}
\end{theorem}

The proof assembles the model-Chern construction, the scalar--curvature
finite trace, and the supported-complex construction. Their proofs,
together with the detailed analytic interfaces, are given below. The present theorem does not assume that
$Z$ is homologically trivial.

\begin{proof}
\textit{Step 1: actual model characters.}
Theorem~\ref{h22c:chern} constructs separately closed model characters,
strictly additive on model vector-bundle exact sequences and functorial
for morphisms of projective flat models. Its proof first establishes
the numerical detector for polynomials in model first Chern forms,
then descends symmetric root polynomials along local sections of a
flag bundle. The actual/CLD comparison and the detector are explained
in Section~\ref{h22c:section}. This construction uses no
finite trace or prescribed-cycle Green pair.

\textit{Step 2: global primitives for generic Chow relations.}
If a model complex is generically acyclic, apply the supported-complex
construction with the entire generic fibre as its exact open. The
relative line construction is then finite flat everywhere on that
fibre, and its scalar--curvature line sum has empty generic divisor.
The trace formula gives a global actual form $\Psi$ whose $D$ is the
alternating model character. Actual-current injectivity turns the
testing identity into an equality of actual forms. This is the
argument of Corollary~\ref{r11p:acyclic}; it also covers the empty
generic support, for which no codimension is assigned to the empty set.

A relation in $K_0(X)_{\mathbf Q}$ is a finite rational combination of
short exact sequence relations and isomorphic-bundle differences.
On a common model, the graph of each generic subbundle yields an exact
model sequence. Strict Whitney reduces the original relation to
differences of model extensions of isomorphic generic bundles.
After uniformizer rescaling, these differences are represented by
generically acyclic two-term complexes. Adding their primitives proves
Proposition~\ref{r11p:generic-k-primitive}.

Now suppose that a model virtual class $\zeta$ satisfies only
$\operatorname{ch}_r(\zeta|_X)=0$. Put
\[
 e_r(T)=\prod_{\substack{0\le j\le n\\j\ne r}}
                  \frac{T-2^j}{2^r-2^j},\qquad
                         \zeta_r=e_r(\psi^2)\zeta.
\]
The rational Chern character is an isomorphism
$K_0(X)_{\mathbf Q}\simeq\bigoplus_j\operatorname{CH}^j(X)_{\mathbf Q}$.
Consequently $\zeta_r|_X=0$. The strict model Adams identity gives
$\operatorname{ch}_r(\bar\zeta_r)=\operatorname{ch}_r(\bar\zeta)$.
The preceding paragraph therefore supplies a global actual primitive
of $\operatorname{ch}_r(\bar\zeta)$. This is
Proposition~\ref{r11p:generic-chow-primitive}. All coefficients are
rational; the construction imposes no restriction on the residue
characteristic.

\textit{Step 3: the Chow map and its pullbacks.}
For $a\in\operatorname{CH}^r(X)_{\mathbf Q}$ choose a generic virtual
bundle with degree-$r$ character $a$, and extend its finitely many
bundles to a common projective flat model. Define
$\mathfrak o^r_{X,K}(a)=[\operatorname{ch}_r(\bar\xi)]$.
Two choices can be pulled to a common model refinement, preserving
their metrics. Their difference has zero generic degree-$r$ Chow
character. Step 2 shows that their actual characters differ by $D$
of a global actual form. Thus the definition is independent of all
choices and is rational-linear.

For $f:X\to Y$, take the closure of its graph in projective models
and refine the source model. The target model bundles now have model
pullbacks. Strict model-Chern functoriality identifies their actual
characters with the genuine analytic pullbacks. The latter commute
with $D$ by the actual pullback theorem, so the equality descends to
the displayed quotient. On Chow groups this is the usual
contravariant operational pullback on smooth varieties.

\textit{Step 4: the prescribed cycle and its exact complement.}
Write $Z=\sum_j a_j[Z_j]$ with distinct integral components and
$a_j\ne0$. Resolve each $\mathcal O_{Z_j}$ by generic vector bundles
and extend these complexes to a common projective flat model.
The degree-$r$ localized character is exactly $[Z_j]$ on its support.
Theorem~\ref{r11p:complex-green} gives an actual Green pair for each
complex on $X^{\mathrm{an}}\setminus Z_j^{\mathrm{an}}$.
Lemma~\ref{r12i:integral-PC} supplies this construction for smooth
integral $X$, including the case that $X$ is not geometrically integral.
The lemma works with native $K$-schemes throughout.

Sum the resulting pairs with coefficients $a_j$. The common domain
is precisely $U_Z$, and the localized characters add to the original
cycle with its original coefficients. The curvature is the degree-$r$
model character of the corresponding rational virtual complex, whose
generic character is $[Z]$; Step 3 therefore identifies its class with
$\mathfrak o^r_{X,K}([Z])$.

For any displayed test $\eta$, the support of $D\eta$ is compact and
disjoint from $|Z|^{\mathrm{an}}$, since its bidegree is
$(n-r+1,n-r+1)$ and $\dim Z=n-r$. This is the support-avoidance
corollary in the actual calculus. Thus its pairing with $G_Z$ is
defined on the exact open. The global curvature term and the cycle
term are also actual top-degree integrals. The all-actual identity
from the supported-complex theorem gives~\eqref{h22:curved-identity}.
For $Z=0$, choose both forms to be zero.
\end{proof}

\begin{corollary}[Curvature correction]
\label{h22:correction}
If $\mathfrak o^r_{X,K}([Z])=0$, the prescribed cycle $Z$ has an actual
zero-curvature Green form on $U_Z$ against all compact actual tests.
\end{corollary}
\begin{proof}
The definition of the quotient supplies a global actual $\Psi_Z$ with
$D\Psi_Z=\Omega_Z$. Set $\gamma_Z=G_Z-\Psi_Z|_{U_Z}$. Actual Stokes
subtracts the curvature term in~\eqref{h22:curved-identity}, giving
\[
          \int_{U_Z}\gamma_Z\wedge D\eta=-\int_{Z^{\mathrm{an}}}\eta.
\]
The correction is global and creates no new deleted subset.
\end{proof}

The additional vanishing
$\mathfrak o^r_{X,K}(\operatorname{CH}^r(X)_{\mathbf Q,\mathrm{hom}})=0$
is the remaining general source-vanishing problem on this route.
The theorem above constructs the map and its curved representatives;
it does not assert this further vanishing.


\section{Model Chern forms: the measure comparison and descent interfaces}
\label{h22c:section}

Fix a finite extension $K/\mathbf Q_p$ with its given absolute value, and
write $R=K^\circ$.  Let $\mathcal X$ be an integral projective flat
$R$-scheme with smooth integral projective generic fibre $X$ of dimension
$n$.  Write $D=d'd''$, and use Mihatsch's actual sheaf $\mathcal B$ on the
ordinary analytic topology.  In its intrinsic trigrading,
$\mathcal B^{0,0,r}$ has bidegree $(r,r)$.  All current comparisons in this
section may be restricted to the original smooth CLD test forms.

The analytic inputs used below are the following primary-source results.
Mihatsch, \emph{$\delta$-Forms on Lubin--Tate spaces},
arXiv:2112.10018v4, Definition 4.2, Remark 4.3(1), Lemma 4.4,
Theorem 4.5, Definition 4.10 and Proposition 4.11 provide the actual
sheaf, pullback, coefficientwise multiplication by PWS forms, integration,
Stokes, and the injective map to currents commuting with $d',d''$.
The paragraph following his Theorem 5.3 identifies model first Chern
forms with elements of $\mathcal B^{0,0,1}$.
Gubler--Martin, \emph{On Zhang's semipositive metrics}, Theorem 1.1 and
Proposition 2.15 give algebraization of model metrics on a proper
variety and uniform density of rational model functions.
Chambert-Loir--Ducros, arXiv:1204.6277v2, Corollary 17.8.9,
\S8.8.5 and Theorem 18.5.5 give the difference-of-psh decomposition,
the multilinear Bedford--Taylor product, and its native model mass
formula.  Proposition 18.5.8 gives the component multiplicities on a
relatively normal model over the discretely valued field $K$.
These references are numbered as in the PDF, rather than the flattened
section counters of its HTML rendering.

\begin{lemma}[Comparison of actual model products and Bedford--Taylor products]
\label{h22c:comparison}
For any finite list of model metrized algebraic line bundles on $X$, the
current of the actual wedge of their first Chern forms equals their
CLD multilinear curvature product, when tested against original smooth
forms.  In top degree the two sides therefore define the same signed
Radon measure.  This assertion holds over $K$ itself.
\end{lemma}

\begin{proof}
Work on an ordinary open on which the bundles are trivial, and write
$u_i=-\log\|s_i\|$ for their continuous PWS weights.  By Corollary
17.8.9 of CLD, model metrized line bundles on a projective
analytification are differences of psh metrized line bundles.  The
identity can consequently be proved by induction for psh factors and
extended by multilinearity.  The local weights and all products below
are actual forms; in particular $Du_i$ is the actual first curvature.

The induction step requires a multiplication statement, which we make
explicit.  Let $T$ be the actual product of the already treated
curvatures.  Its current has order zero: on a common moment
presentation, pairing $T$ with a fixed compactly supported smooth
complementary form gives a finite sum of signed measures with smooth
coefficient densities on polyhedra.  For a continuous PWS weight $u$,
Mihatsch's coefficientwise formula in Remark 4.3(1) gives
\[
             [uT](\eta)=\int u\,d\mu_{T,\eta},
       \qquad \mu_{T,\eta}\text{ the measure of }T\wedge\eta .
\]
The formula can be checked on a finite ordinary cover of the support
of $\eta$ and then summed using smooth partitions.  Thus it is
independent of the common moment presentation.  By the induction
hypothesis, $[T]$ is the CLD curvature product.  Its standard
continuous multiplication by $u$ gives precisely the same coefficient
measures.  Equivalently, equality on smooth scalar multipliers extends
to continuous scalar multipliers by uniform approximation on each
compact support.  Hence $[uT]=u[T]$.

The factors of $T$ are separately closed.  Actual Leibniz and the
chain-map assertion of Mihatsch's Proposition 4.11 now give
\[
 [Du\wedge T]=[D(uT)]=D[uT]=D\bigl(u[T]\bigr).
\]
This is the defining Bedford--Taylor recursion, with its extension to
differences of psh functions supplied by CLD \S8.8.5.  The induction
starts with the fundamental integration current and proves the
assertion in every degree.  In top degree, these currents are Radon
measures and equality on smooth scalar tests determines the measure.
No scalar extension or general product of arbitrary currents occurs.
\end{proof}

\begin{lemma}[Detection of a closed pure actual form]
\label{h22c:detector}
Let $0\le r\le n$, let $q=n-r$, and let
$\Omega\in\mathcal B^{0,0,r}(X^{\mathrm{an}})$ be separately closed.
Suppose that
\[
                 \Omega\wedge D\phi_1\wedge\cdots\wedge D\phi_q=0
\]
as a signed Radon measure for every tuple of global rational model
functions $\phi_1,\ldots,\phi_q$.  Then $\Omega=0$ as an actual form.
\end{lemma}

\begin{proof}
A compactly supported actual top form has a finite signed measure by
its local polyhedral integration formula.  Its associated current is
integration against that measure.  Thus the case $q=0$ follows from
the injectivity in Mihatsch's Proposition 4.11.

For $q>0$, successively replace model-function slots by global
original-smooth functions.  At an induction step, let $T$ denote
$\Omega$ wedged with the other $q-1$ Hessians.  It is separately
closed and has bidegree $(n-1,n-1)$.  If $\phi$ is a global rational
model function and $\psi$ is a global original-smooth function,
actual Stokes on the compact space $X^{\mathrm{an}}$ gives
\[
                  \int\phi\,T\wedge D\psi
                    =\int\psi\,T\wedge D\phi=0.
\]
The last equality is the induction hypothesis.  The expression on
the left tests the fixed finite signed measure $T\wedge D\psi$
against $\phi$.  Uniform density of rational model functions in
$C(X^{\mathrm{an}})$ therefore makes that measure zero.  After $q$
steps,
\[
                  \Omega\wedge D\chi_1\wedge\cdots\wedge D\chi_q=0
\tag{H22.C1}\label{h22c:all-smooth-hessians}
\]
for all global original-smooth functions $\chi_i$.  This argument
does not differentiate a uniformly convergent sequence.

We spell out the dimension of a witnessing skeleton.  By Mihatsch's
Definition 3.8, $\Sigma(U,g)$ is the $n$-dimensional locus of the
moment skeleton; by his Proposition 3.7 it is purely
$n$-dimensional away from the analytic boundary.  His Lemma 4.4,
applied on an ordinary open and with refining toric coordinates,
therefore implies that a nonzero actual form has a nonzero restriction
on such an $n$-dimensional skeleton.  This can also be seen by first
using Proposition 4.11 to choose a detecting smooth complementary
test, whose top-degree product has a containing $n$-dimensional
skeleton by Corollary 4.9 and Definition 4.10.

Suppose $\Omega\ne0$ and choose such a restriction.  After a
polyhedral subdivision, purity means that on some $q$-dimensional
open face $\sigma$ its coefficient is a nonzero scalar function.
The moment coordinates $u_a=-\log|g_a|$ have an injective tangent
map on each linearity face: a nonzero tangent kernel would give
positive-dimensional fibres, contrary to the finite-fibre property
of skeleton tropicalization.  Hence
\[
                         h=\frac12\sum_a u_a^2
\]
has positive definite Hessian on $T\sigma$.  Choose a smooth cutoff
supported inside the moment-chart open and equal to one near a point
where the scalar coefficient is nonzero.  Multiplying $h$ by this
cutoff and extending by zero gives a global original-smooth function
$\chi$ with $D\chi=Dh$ on that neighbourhood.  The coefficientwise
wedge formula makes
$\Omega\wedge(D\chi)^q$ nonzero on $\sigma$.
This contradicts \eqref{h22c:all-smooth-hessians} and top-form
injectivity.  Therefore $\Omega=0$.
\end{proof}

\begin{theorem}[Numerically trivial model-root polynomials vanish as actual forms]
\label{h22c:numerical-zero}
Let $P$ be a homogeneous rational polynomial of degree $r$ in the first
Chern classes of finitely many model line bundles $\mathcal L_i$ on
$\mathcal X$, where $1\le r\le n$.  If
\[
 \xi=P\bigl(c_1(\mathcal L_{i,s})\bigr)\in A^r(\mathcal X_s)_{\mathbf Q}
\]
has degree zero on every integral $r$-dimensional subvariety of
$\mathcal X_s$, then
\[
                    P\bigl(c_1(\overline L_i)\bigr)=0
                    \quad\text{in }\mathcal B^{0,0,r}(X^{\mathrm{an}}).
\]
\end{theorem}

\begin{proof}
Put $\Omega=P(c_1(\overline L_i))$ and $q=n-r$.  For any $q$ rational
model functions, take a common projective flat model refinement on
which they are represented by vertical rational Cartier divisors
$D_1,\ldots,D_q$.  We may moreover normalize this refinement.
Indeed $R$ is an excellent DVR, so normalization is finite; it
preserves projectivity and, because $X$ is smooth, is the identity
on the generic fibre.  The resulting integral model remains flat.
Write $f:\mathcal Y\to\mathcal X$ for the resulting map.  As usual,
common models of finitely many model metrics are obtained by
admissible modifications; their algebraization in the projective
case is the model-metric correspondence used in Gubler--Martin
Theorem 1.1 and in the proof of CLD Corollary 17.8.9.

Lemma~\ref{h22c:comparison} identifies the top actual wedge with the
CLD model measure.  On the normal model, CLD Theorem 18.5.5 and
Proposition 18.5.8 identify its mass at the divisorial point of an
irreducible special-fibre component $V$ with
\[
 b_V\,\deg\!\left(
     f_s^*\xi\cap
     \prod_{j=1}^q c_1(\mathcal O(D_j)|_V)\cap[V]\right),
\]
where $b_V$ is the full special-fibre multiplicity.  Rational divisor
denominators are retained multilinearly.  Alternatively, without
normalization one may keep the full native weight
$m(x)\operatorname{def}_K(x)\lambda_{\mathcal Y}(x)$ from
Theorem 18.5.5; the following zero calculation is unchanged.

The operational projection formula rewrites the degree as
\[
 \deg\!\left(\xi\cap f_{s*}
       \left(\prod_{j=1}^q c_1(\mathcal O(D_j)|_V)\cap[V]\right)\right).
\]
The inner intersection is a rational $r$-cycle.  Proper pushforward
is a rational combination of integral $r$-cycles on $\mathcal X_s$;
components whose images have smaller dimension contribute zero.
The numerical hypothesis therefore makes every displayed degree
zero.  The support assertion in CLD Theorem 18.5.5 now makes the
entire top measure zero.  Lemma~\ref{h22c:detector} gives $\Omega=0$.
Neither smoothness nor reducedness of the special fibre is used.
\end{proof}

\begin{theorem}[Strict actual Chern forms of model vector bundles]
\label{h22c:chern}
For model vector bundles $\mathcal E_1,\ldots,\mathcal E_s$ on
$\mathcal X$ and a homogeneous rational Chern polynomial $P$ of
degree $r$, $1\le r\le n$, there is a unique separately closed
actual form
\[
                c_{\mathcal X}(P)\in\mathcal B^{0,0,r}(X^{\mathrm{an}})
\]
whose pullback to the simultaneous full flag bundle is the
corresponding model-root polynomial.  The construction is functorial
under morphisms of integral projective flat models with smooth
integral generic fibres.  Polynomial operational identities, including
Whitney identities, Chern-character additivity and Adams identities,
hold as strict identities of actual forms.  Numerical triviality
of the special-fibre polynomial implies $c_{\mathcal X}(P)=0$.
\end{theorem}

\begin{proof}
Let
$a:\mathcal F=\prod_{\mathcal X}\operatorname{Flag}(\mathcal E_j)
\to\mathcal X$ and write $\Omega_F$ for its model-root polynomial.
Flag bundles, and their fibre products here, are integral projective
flat models; their generic fibres are smooth integral varieties.
On $\mathcal F\times_{\mathcal X}\mathcal F$, the two root
polynomials represent the same operational Chern class, by the two
tautological filtrations.  Thus their difference restricts to the
zero operational class on the special fibre.  It is a polynomial
in first Chern classes of model line bundles, so
Theorem~\ref{h22c:numerical-zero}, applied in the dimension of this
larger generic fibre, proves
\[
                         p_1^*\Omega_F=p_2^*\Omega_F
\tag{H22.C2}\label{h22c:flag-cocycle}
\]
as actual forms.

Trivialize the finitely many generic bundles simultaneously on
Zariski opens $U_i\subset X$.  Their standard flags give algebraic
sections $s_i:U_i\to F$.  These are genuine analytic sections on
ordinary opens after analytification.  All their domains and targets
are pure-dimensional and boundaryless, so Mihatsch's Theorem 4.5
applies even though the sections need not extend to the models.
Pulling \eqref{h22c:flag-cocycle} back by $(s_i,s_j)$ proves that
$s_i^*\Omega_F$ glue to an actual form $c_{\mathcal X}(P)$.
Pulling it back by $(\mathrm{id},s_i\circ a)$ on $a^{-1}(U_i)$
proves $a^*c_{\mathcal X}(P)=\Omega_F$.  The sections also make
$a^*$ injective, proving uniqueness.  Pullback respects tridegree
and differentials, so the descended form is pure and separately
closed.

To prove a universal operational identity among finitely many
bundles, take their simultaneous flag bundle and subtract the
corresponding root polynomials.  The difference is zero in
special-fibre operational Chow.  The preceding theorem makes it
zero as an actual form upstairs; injectivity of flag pullback gives
the strict identity downstairs.  This proves Whitney and additive
relations.  Applied to the universal symmetric-root identity for
Adams operations, it gives
\[
          \operatorname{ch}_r(\psi^k\xi)=k^r\operatorname{ch}_r(\xi)
\]
strictly for model virtual bundles and $k\ge1$.

For a morphism of the specified models, pull back the flag bundle
and its tautological line quotients.  Model first Chern forms are
functorial, so the pulled-back form and the newly constructed form
have the same characteristic flag pullback.  Faithfulness of that
pullback proves functoriality.  The same argument with the product
of two splitting presentations proves independence of presentation.
In particular a model refinement leaves the actual form of a fixed
pulled-back model bundle unchanged.

Finally, if $\xi=P(c(\mathcal E_{j,s}))$ is numerically trivial,
then for every integral $r$-cycle $W$ on $\mathcal F_s$,
\[
             \deg(a_s^*\xi\cap[W])
               =\deg(\xi\cap a_{s*}[W])=0.
\]
Theorem~\ref{h22c:numerical-zero} kills $\Omega_F$ and the local
sections kill its descent.  This proves the last assertion.
\end{proof}

\paragraph*{Scope of the conclusion.}
The preceding construction uses no finite trace and no prescribed
Green-pair theorem.  It supplies the strict model-Chern input used in
Sections~\ref{r11p:minimal-trace} and~\ref{r11p:supported-complex}.  It does not identify two arbitrary model extensions merely
because their generic bundles are isomorphic; that comparison requires
the separate global-primitive argument.  It also makes no assertion
that the curvature of every geometrically homologically trivial
cycle is exact.


\section{The minimal trace needed by a supported perfect complex}
\label{r11p:minimal-trace}

\paragraph*{The scalar--curvature module.}
The supported-complex Green formula needs only expressions
\[
                       u_0\,Du_1\cdots Du_k, \qquad D=d'd'',
\tag{MT1}\label{r11p:module-generators}
\]
with PL scalar potentials. A single parameter evaluates the whole
scalar coefficient $u_0$, even when $u_0$ is nonlinear PL in many moment
coordinates. The proof constructs this trace directly from scalar PL
descent, max/min sorting, and Euclidean projection. The inputs are
specified below; a trace on arbitrary actual forms or a universal
Green-pair construction is not assumed.

Throughout this section $K$ is a complete nontrivially valued field of
characteristic zero, and all analytic spaces are the analytifications of
normal integral separated $K$-varieties, with their ordinary topology.
All forms are Mihatsch actual delta-forms.  A PL scalar means a function
locally presented by a finite piecewise affine function of logarithms of
analytic units.  Its $D$ is a pure form of tridegree $(0,0,1)$.
Let $\mathscr M^k$ be the actual image sheaf of finite sums of
\eqref{r11p:module-generators}.  Relations in $\mathscr M^k$ are actual-form
relations; no chosen presentation is part of a section.

\subsection{Primary analytic inputs and native degree}

We use the following statements in their stated scopes.
\begin{enumerate}
\item Mihatsch, \emph{Delta-Forms on Lubin--Tate Spaces},
arXiv:2112.10018v4, Definition 4.2 and Theorem 4.5: local Euclidean
presentations, the ordinary sheaf, arbitrary actual pullback between
pure-dimensional boundaryless spaces, and compatibility with products
and differential operators.
\item Ibid., Remark 4.3(1), Lemma 4.4, Corollary 4.9 and Proposition 4.11:
coefficientwise multiplication by a PWS scalar on a witnessing skeleton;
support avoidance when the maximum bidegree exceeds the algebraic
dimension; Stokes; and injectivity into ordinary CLD currents.
\item Ibid., Definition 3.4 and Lemmas 3.3 and 3.15:
native calibrated tropical weights and the finite-flat degree formula.
These are used for top forms only.
\item Ibid., Proposition 2.17, Definition 2.29(3), Lemma 2.37,
Proposition 2.38 and Definition 2.41: proper-on-support pushforward
of polyhedral currents on weighted pwl spaces, its differential
compatibility, and the boundary-excluded chart characterization of
tropical pullback.  Definition 3.8, Remark 3.9 and the coordinate case
of the proof of Proposition 3.21 identify product skeletons with their
linear functions and native weights.  Definition 4.10 defines analytic
integration intrinsically on skeletons.
\item Mihatsch, \emph{On Tropical Intersection Theory},
arXiv:2107.12067v3, Section 2.3, Example 2.5, Theorem 4.1(3), and
Proposition/Definition 4.2: affine proper-on-support Euclidean
pushforward, its projection formula, and affine pullback.
\end{enumerate}

For a finite dominant map $q:Y\to X$ of generic degree $d$ and any compact
actual top form $\tau$ on $X$, these inputs give
\[
                            \int_Yq^*\tau=d\int_X\tau.
\tag{MT2}\label{r11p:native-degree}
\]
Indeed remove a proper closed subset outside which $q$ is finite flat of
constant degree $d$.  The support of $\tau$ avoids this subset by
Corollary 4.9; so does the support of its pullback above it.
Lemma 3.15 and the native skeletal integration formula apply on the
remaining open.  This argument does not construct a trace of arbitrary
source coefficients.  It also proves that finite dominant pullback on
actual forms is faithful: pair a putatively zero pullback with pulled-back
original-smooth complementary tests, use \eqref{r11p:native-degree}, and
apply Proposition 4.11.

\subsection{Only scalar PL functions pass through a finite quotient}

\begin{lemma}[Ordinary local scalar PL descent]
\label{r11p:pl-descent}
Let $q:Y\to X$ be finite Galois between normal integral varieties, with
group $G$.  A $G$-invariant source PL scalar descends to a PL scalar on
$X$.  If finitely many invariant potentials on
$Y^{\rm an}\times(\mathbf G_m)^{\rm an}$ are jointly PL in fixed source
unit logarithms and the parameter $t=v(z)$, their descended potentials
are locally jointly PL in a fixed finite tuple of base-unit logarithms
and $t$.  The base tuple and base shrinking are independent of $t$.
\end{lemma}

\begin{proof}
Work around a whole finite analytic fiber.  Choose disjoint ordinary
neighborhoods of its points on which the finitely many source potentials
have moment presentations.  Properness permits shrinking the base so
that its full inverse image lies in their union.  The corresponding
source pieces are then clopen.  Analytic units can be pasted on these
pieces.  Add constant units with distinct valuations as markers; after
shrinking compact neighborhoods, their moment images lie in separated
compact polyhedral boxes.  Finite PL extension on these boxes
combines the presentations.  A smooth cutoff of a PL potential is
neither required nor used.

Include the full $G$-orbit of the resulting source-unit tuple.  Its
logarithms $\ell\in\mathbf R^s$ then carry a coordinate-permutation
$G$-action; add a constant unit if necessary so that $s\geq1$.
If an ambient PL presentation $\varphi(\ell)$ is
invariant only after restriction to the actual source moment image,
replace it by $|G|^{-1}\sum_{g\in G}\varphi(g\ell)$.  This is an
ambient invariant PL function with the same source restriction.
For ordered orbit potentials the sorting itself already gives this
ambient invariance.  The same operation leaves an independent
parameter fixed.  Put $d=|G|$ and $N=d(s-1)+1$.  For distinct integers $a_j$
take the Vandermonde directions
\[
                     m_j=(1,a_j,\ldots,a_j^{s-1})\in\mathbf Z^s.
\]
Define $J(\ell)$ by recording, for each $j$, the ordered multiset
$\{\langle m_j,g\ell\rangle:g\in G\}$.  It separates $G$-orbits.
In fact if $J(x)=J(y)$, for each $j$ one can choose $g_j$ with
$\langle m_j,x\rangle=\langle m_j,g_jy\rangle$.  One $g$ occurs for
at least $s$ indices.  The resulting polynomial of degree at most $s-1$
has $s$ roots, hence $x=gy$.

The map $J$ is proper: any $s$ independent Vandermonde directions bound
$\ell$ when all the displayed coordinates are bounded.  It is finite
piecewise linear.  On a common polyhedral subdivision its linearity
cells map affinely and injectively to their images; a noninjective
tangent direction would produce an infinite fiber.  Every invariant
PL scalar therefore factors as a PL scalar through $J$, after common
image subdivision and PL ambient extension.  The identical argument
applies to $J\times\operatorname{id}_{\mathbf R}$ and jointly PL
potentials with the fixed parameter $t$.

It remains to realize $J$ downstairs without an illicit nonlinear
pullback of forms.  For each source monomial $F_j=f^{m_j}$ take its monic
orbit polynomial
\[
                          \prod_{g\in G}(T-gF_j).
\]
Its coefficients are invariant analytic functions, hence analytic
functions on the normal finite quotient.  Algebraically the quotient
is $X$; analytically this follows on affinoids from finite coherent
base change and the Reynolds idempotent in characteristic zero.
With the convention that the coefficient of $T^i$ gives the point
$(i,v(a_i))$, the ordered root valuations are the negatives of the
lower Newton polygon's slopes, with multiplicity.  They are therefore
PL in the valuations of the coefficients that contribute to that
polygon.

At a fixed base point retain the nonzero coefficients and shrink so
that they are units.  The constant coefficient $a_0$ is a unit because
it is a product of source units; the leading coefficient is one.
Choose a uniform upper bound $B$ for $v(a_0)$ on the shrunken base.
Every coefficient vanishing at the fixed point can, by continuity and
a further common shrinking, be made to have valuation greater than
$\max(B,0)+1$.  Its point is then above the endpoint chord of the
Newton polygon and cannot affect the lower hull.  Such coefficients
can be omitted uniformly.
Thus $J$ is PL in finitely many base-unit logarithms.  These
coefficients involve source units alone, so their choice and shrinking
are independent of $t$.  Substitution is only a substitution into a
scalar PL function.  Taking its actual derivatives happens afterwards;
no nonlinear PL pullback of a general delta-form has been postulated.
\end{proof}

\subsection{Sorting gives a pure-curvature trace}

\begin{lemma}[Max/min identities]
\label{r11p:sorting}
For source PL scalars $u,v$, put $M=\max(u,v)$ and $m=\min(u,v)$.
For every $k\geq0$ there are actual identities
\[
 (DM)^k+(Dm)^k=(Du)^k+(Dv)^k,
\tag{MT3}\label{r11p:sorting-powers}
\]
\[
 M(DM)^k+m(Dm)^k=u(Du)^k+v(Dv)^k.
\tag{MT4}\label{r11p:sorting-energies}
\]
\end{lemma}

\begin{proof}
Write $h=u-v$, $h_+=\max(h,0)$, $h_-=\min(h,0)$, and
$A=Dh_+$, $B=Dh_-$.  Coefficientwise multiplication on support
polyhedra gives $h_+B=h_-A=0$: on one open side the appropriate
function is zero, and on the common wall its continuous scalar
coefficient is zero.  Applying $D$ and using separate closedness gives
$AB=0$.  Since
\[
 M=v+h_+,\quad m=v+h_-,\quad u=v+h_++h_-,
\]
expansion proves \eqref{r11p:sorting-powers}.  The same expansion,
now using both $AB=0$ and the two scalar-weighted zero identities,
proves \eqref{r11p:sorting-energies}.
\end{proof}

The second identity is useful but has a definite limitation:
polarizing $u(Du)^k$ recovers symmetrized energies, not an arbitrary
ordered scalar-curvature monomial $u_0Du_1\cdots Du_k$.
The parameter argument below supplies precisely that missing step.

\begin{lemma}[Trace of pure PL curvature polynomials]
\label{r11p:pure-trace}
For finite dominant $q:Y\to X$ between normal integral varieties, a
finite sum $C$ of products $Du_1\cdots Du_k$ has an intrinsic actual
trace $\operatorname{Tr}^{\rm curv}_q(C)$, characterized by
\[
 \int_X\operatorname{Tr}^{\rm curv}_q(C)\wedge\eta
                         =\int_YC\wedge q^*\eta
\tag{MT5}\label{r11p:pure-pairing}
\]
for all compact complementary actual tests.  With a parameter, jointly
PL source potentials in fixed source logarithms and $t$ yield a trace
with a product-moment presentation using fixed base logarithms and $t$.
\end{lemma}

\begin{proof}
First let $q$ be Galois.  Sort the orbit of a scalar $u$ by a finite
sorting network.  Repeated use of \eqref{r11p:sorting-powers} gives
\[
                     \sum_{g\in G}(Dg^*u)^k
                            =\sum_{j=1}^d(Du_{(j)})^k.
\]
The ordered potentials $u_{(j)}$ are invariant and descend by
Lemma~\ref{r11p:pl-descent}.  The right side is consequently the
pullback of a downstairs actual form.  Define it to be the trace of
$(Du)^k$.  Polarization,
\[
 Du_1\cdots Du_k
   =\frac1{k!}\sum_{I\subset\{1,\ldots,k\}}
           (-1)^{k-|I|}\left(D\sum_{i\in I}u_i\right)^k,
\]
does the same for all mixed products.  For $k=0$ set
$\operatorname{Tr}^{\rm curv}_q(1)=d$.

This construction satisfies the actual orbit identity
\[
                 q^*\operatorname{Tr}^{\rm curv}_q(C)
                                      =\sum_{g\in G}g^*C.
\]
Wedge with $q^*\eta$, integrate, and use
\eqref{r11p:native-degree} and invariance of native integration under
automorphisms.  Division by $d$ proves \eqref{r11p:pure-pairing}.
Finite pullback faithfulness proves independence of every actual
relation.  Current injectivity gives gluing on ordinary overlaps.
The pairing permits arbitrary actual downstairs multipliers; they
have not been treated as arbitrary source coefficients.

For a general separable finite map take a normal Galois closure
$\widetilde Y\to X$ factoring through $r:\widetilde Y\to Y$, and divide
the Galois trace of $r^*C$ by $e=\deg r$.  Applied to the top form
$Cq^*\eta$, \eqref{r11p:native-degree} proves the required pairing and
independence of the closure.  Characteristic zero covers every map
used here.  All spaces and parameter tori remain over the original
$K$; native generic degrees are not normalized away.

For the last assertion use the parameter version of
Lemma~\ref{r11p:pl-descent} at the sorting step.  The descended
potentials have fixed base moments, so their Hessians and products
have the required product-moment presentation.
\end{proof}

\subsection{Certified integration along a parameter torus}

\paragraph*{The boundary issue.}
Compact analytic support does not imply that an ambient moment
coefficient avoids the tropicalized analytic boundary of an affinoid
chart.  We do not use the single-chart integral formula with that
unverified hypothesis.  Instead we first prove a pushforward identity
for polyhedral currents on the product skeleton, retaining the boundary
exclusion of Proposition 2.38 at every chart calculation.  The
intrinsic skeletal definition of integration then gives the analytic
pairing.  In this subsection $f^\star$ denotes tropical-space pullback
in that proposition, distinguished from the notation $t_c^*$ for the
resulting actual analytic form.

\begin{lemma}[The coordinate product skeleton, including its structures]
\label{r11p:coordinate-product-skeleton}
Let $V$ be a boundaryless ordinary analytic open of pure dimension $n$,
let $b$ be a tuple of analytic units on $V$, and put $S=\Sigma(V,b)$.
The relative Gauss section identifies
\[
 \begin{split}
 \iota_b:S\times\mathbf R&\ \xrightarrow{\ \sim\ }\
       \Sigma(V\times(\mathbf G_m)^{\rm an},(b,z)),\\
                  (y,t)&\longmapsto \nu_t(y),
 \end{split}
\tag{PS1}\label{r11p:product-skeleton-identity}
\]
where $\nu_t(y)$ is the Gauss point of radius $e^{-t}$ in the torus
over $\mathcal H(y)$.  This is an isomorphism of weighted pwl spaces
and of tropical spaces.  Its weight is the product of the native
weight on $S$ and the standard weight on $\mathbf R$, and its linear
function sheaf is generated by the base linear functions and $t$.
The identification is compatible with every ordinary base shrinking
and every additional base moment refinement.
\end{lemma}

\begin{proof}
Let $x$ lie above $y$ in the analytic product.  The graded field
generated by $b$ and $z$ has transcendence degree $n+1$ precisely
when the graded field generated by $b$ at $y$ has transcendence
degree $n$ and $\widetilde z$ is transcendental over
$\widetilde{\mathcal H(y)}$.  Indeed adjoining $z$ increases that
degree by at most one; in the degree-$n$ case
$\widetilde{\mathcal H(y)}$ is algebraic over the graded field
generated by $b$.  This is the graded-transcendence calculation in
(3.27) of the cited proof, here with the single coordinate $z$.

The condition on $\widetilde z$ characterizes the Gauss valuation
over $\mathcal H(y)$: no cancellation among leading monomials is
possible, and for every Laurent polynomial one has
\[
 \left|\sum_i a_i z^i\right|_{\nu_t(y)}
                         =\max_i |a_i|e^{-it}.
\]
Conversely these Gauss valuations have the required graded
transcendence degree.  Thus Definition 3.8 and the absence of analytic
boundary give the claimed bijection.  No valued extension of $K$,
or rationality of the radius in its value group, is needed.

The inverse is $(p,v(z))$ on the displayed skeleton.  It is pwl:
a base skeleton chart is described by an affinoid domain and
valuations of its analytic units; their pullbacks are valuations of
the same units on the product, and $v(z)$ is the extra coordinate.
The corresponding closed chart domains pull back to closed pwl
domains.  This inverse is also proper.  A compact subset of
$S\times\mathbf R$ is contained in a finite union of compact base
analytic domains times a bounded interval.  Its inverse image is
closed in the corresponding finite union of products with compact
analytic annuli, and is therefore compact.  It follows that the
continuous bijective inverse is a homeomorphism.  Subdivision into
local polyhedral charts makes its injective affine pieces invertible
onto their images, so both directions are pwl.  This is the
coordinate specialization of the topology and atlas argument
following (3.31), now without its changes of field or center.

For the weights, take a top-dimensional base cell and a local toric
projection used in Definition 3.8(2) to compute its native degree
$e$.  Its product with the identity on $\mathbf G_m$ is finite flat
of the same degree $e$.  Hence its induced weight is the base weight
times the standard parameter weight.  This calculation retains
residue-field degrees and is exactly the last weight calculation
in the proof of Proposition 3.21.  It is local on cells and hence
determines the entire product weight.

Under the bijection the moments are $(t_b,t)$.
Definition 3.8(3) therefore gives precisely the product linear
function sheaf, not merely a set-theoretic identification.
For additional base moments $b'$, the same Gauss section applies.
The inclusions of base and product skeletons agree under it, and
their weights agree by Remark 3.9.  The new linear functions are the
new base ones together with the same $t$.  This proves the stated
compatibility for all refinements and ordinary restrictions.
\end{proof}

\begin{lemma}[Tropical product pushforward with boundary control]
\label{r11p:tropical-product-push}
Let $S$ be a pure-dimensional tropical space and
$f:S\to\mathbf R^N$ a linear map.  Give $S\times\mathbf R$ the
product weight and linear functions.  Write $p_S$ for its projection
and $\pi:\mathbf R^N\times\mathbf R\to\mathbf R^N$ for the coordinate
projection.  If $\Gamma$ is a Euclidean delta-form with support in
a fixed compact set of parameter values, then
\[
 (p_S)_*\bigl((f,\operatorname{id})^\star\Gamma\bigr)
                                      =f^\star(\pi_*\Gamma)
\tag{PS2}\label{r11p:tropical-product-push-identity}
\]
as polyhedral currents on $S$.  The left side uses only the
proper-on-support pwl-current pushforward of Proposition 2.17.
\end{lemma}

\begin{proof}
The product is a tropical space: on a product polyhedral subdivision,
the vertical balancing conditions are those of $S$, and horizontal
internal facets cancel in pairs with the standard parameter weight.
The support of
$B=(f,\operatorname{id})^\star\Gamma$ is contained in
$S\times J_0$ for a compact interval $J_0$.  Thus it is proper over
$S$, and Proposition 2.17 supplies the polyhedral current
$P=(p_S)_*B$ before any assertion that it is a delta-form.

Work in a polyhedral neighborhood carrying a finite-fiber linear
map $h$ to $\mathbf R^m$, as supplied by the axioms of a sheaf of
linear functions.
Append $f$ and put $g=(f,h)$.  Then $g$ has finite fibers and
$f=A\circ g$, where $A$ is a literal coordinate projection.
Let $H$ be any compact polyhedral subset of the interior of this
neighborhood, and put
\[
 E_H=g(\partial_S H),\qquad
 W_H=\mathbf R^{N+m}\setminus E_H,\qquad
 H_g^\circ=H\setminus g^{-1}(E_H).
\]
The restricted $g:H_g^\circ\to W_H$ is proper.  Write
$C_H=g_*[H_{\dim S},\mu_S]|_{W_H}$.
It is a tropical cycle on $W_H$ by Corollary 2.35.

Choose an interval $J=[a,b]$ containing $J_0$ in its interior.
The boundary of $H\times J$ is contained in
\[
        (\partial_S H\times J)\ \cup\ (H\times\{a,b\}).
\]
Apply the characterizing identity (2.40) of Proposition 2.38 to
this compact product and the refinement $g\times\operatorname{id}$.
On the boundary-free region $W_H\times(a,b)$ it gives
\[
 (g\times\operatorname{id})_*
       \bigl(B|_{H_g^\circ\times(a,b)}\bigr)
       =(C_H\times\mathbf R)\wedge
                             (A\times\operatorname{id})^*\Gamma.
\tag{PS3}\label{r11p:boundary-chart-product}
\]
Both sides vanish near the artificial parameter endpoints because
$\Gamma$ does.  They may therefore be viewed on $W_H\times\mathbf R$,
but the exclusion of $E_H$ in the base remains in force.  There
is no assertion of this chart identity across $g(\partial_S H)$.

Let $\pi_g:W_H\times\mathbf R\to W_H$ be projection.  Functoriality
of pwl-current pushforward and the Euclidean projection formula
give, entirely on $W_H$,
\[
\begin{split}
  g_*\bigl(P|_{H_g^\circ}\bigr)
   &=\pi_{g,*}\!\left((C_H\times\mathbf R)\wedge
                          (A\times\operatorname{id})^*\Gamma\right)\\
   &=C_H\wedge
                  \pi_{g,*}(A\times\operatorname{id})^*\Gamma\\
   &=C_H\wedge A^*(\pi_*\Gamma).
\end{split}
\tag{PS4}\label{r11p:boundary-chart-push}
\]
The last step needs only coordinate Fubini.  For a coordinate
projection $A$, its Euclidean delta-pullback adds free standard
coordinates.  On each weighted carrier of $\Gamma$, fiber
integration in $t$ before or after adding those coordinates
has the same fibers, coefficient integral and quotient weights.
Thus
$\pi_{g,*}(A\times\operatorname{id})^*\Gamma=A^*(\pi_*\Gamma)$.
This carrierwise argument applies to an arbitrary Euclidean
delta-form, with compact parameter support; it does not invoke a
general affine base-change lemma.

The last expression in \eqref{r11p:boundary-chart-push} is exactly
the characterizing expression (2.40) for
$f^\star(\pi_*\Gamma)$.  Varying the compact polyhedral set $H$
inside the fixed finite-fiber $g$-chart, Lemma 2.37 shows that
$P=f^\star(\pi_*\Gamma)$ on that chart.  These charts cover $S$,
which proves \eqref{r11p:tropical-product-push-identity}.
In particular the polyhedral pushforward has the asserted
delta-form representative; this is a conclusion of the calculation.
\end{proof}

\begin{lemma}[Integration on the product-moment subspace]
\label{r11p:product-integration}
Let $p:U\times(\mathbf G_m)^{\rm an}\to U$ be projection.  Suppose
\[
                   \beta=t_{(c,z)}^*\Gamma
\]
for a fixed base analytic-unit tuple $c$, $t=v(z)$, and a Euclidean
delta-form $\Gamma(u,t)$ supported in a fixed compact set of $t$.
With $\pi(u,t)=u$, define
\[
                       I_p(\beta)=t_c^*(\pi_*\Gamma).
\tag{MT6}\label{r11p:product-definition}
\]
This is intrinsic among product-moment presentations and satisfies
\[
             \int_UI_p(\beta)\wedge\eta
                    =\int_{U\times\mathbf G_m}\beta\wedge p^*\eta
\tag{MT7}\label{r11p:product-pairing}
\]
for every compact actual complementary test.  It commutes with
$d',d''$ on these presentations.
\end{lemma}

\begin{proof}
\textit{1. A skeletal identity before integration.}
Euclidean pushforward supplies the actual candidate
$\theta=t_c^*(\pi_*\Gamma)$.
On an arbitrary ordinary $V\subseteq U$, refine the base moments
by tuple concatenation, $b=(c,h)$, and put $S=\Sigma(V,b)$.
Under \eqref{r11p:product-skeleton-identity}, the restriction of
$\beta$ to the product skeleton is
$(t_c|_S,\operatorname{id})^\star\Gamma$.
Lemma~\ref{r11p:tropical-product-push} therefore gives
\[
       \theta|_S
        =(p_S)_*\bigl(\beta|_{S\times\mathbf R}\bigr).
\tag{PS5}\label{r11p:actual-product-skeletal}
\]
This is an equality of intrinsic polyhedral currents on $S$ for
every base refinement.  Its proof used the boundary-excluded chart
identities \eqref{r11p:boundary-chart-product}--\eqref{r11p:boundary-chart-push},
not an integral over a global tropicalization of $V$.

If two product-moment presentations define the same actual $\beta$,
concatenate both base tuples and any further test tuple.  Their
restrictions to every resulting product skeleton agree by Definition
4.2.  Equation \eqref{r11p:actual-product-skeletal} makes their
candidate outputs agree on every corresponding base skeleton.
Lemma 4.4, using all refinements, gives actual equality.  This proves
independence of presentations and their relations, and also their
compatibility on ordinary overlaps, before any current pairing.

\textit{2. Products with arbitrary actual base forms.}
On an open $V$ where an actual form $\eta$ has a moment presentation,
take $b$ to be the concatenation of its moments and $c$, and write
$\eta=t_b^*\eta_b$.  Let $A$ be the coordinate projection selecting
$c$.  The product $\beta\wedge p^*\eta$ has Euclidean coefficient
\[
           (A\times\operatorname{id})^*\Gamma\wedge\pi_b^*\eta_b
\]
in the fixed moments $(b,z)$, with the same compact parameter support.
Euclidean projection and coordinate Fubini give
\[
 \pi_{b,*}\!\left(
      (A\times\operatorname{id})^*\Gamma\wedge\pi_b^*\eta_b\right)
                         = A^*(\pi_*\Gamma)\wedge\eta_b.
\]
Apply \eqref{r11p:actual-product-skeletal} to this coefficient.
For $S=\Sigma(V,b)$ it gives
\[
 (p_S)_*
     \left((\beta\wedge p^*\eta)|_{S\times\mathbf R}\right)
                         =(\theta\wedge\eta)|_S.
\tag{PS6}\label{r11p:product-test-skeletal}
\]
The coefficient $\eta_b$ here is an arbitrary Euclidean delta-form.
No restriction to smooth tests or to products of model curvatures
has occurred.

\textit{3. Intrinsic skeletal integration and ordinary partitions.}
Now let $\eta$ be compact and of complementary bidegree.  Cover its
support by finitely many ordinary opens $V_i$ on which it has the
presentations used in the preceding step, and choose an
original-smooth partition $(\lambda_i)$ subordinate to this cover
on a neighborhood of the support.  Put
$S_i=\Sigma(V_i,b_i)$.
On $V_i$ the unweighted form $\theta\wedge\eta$ is top-degree
and is presented in $b_i$; on $V_i\times\mathbf G_m$ the
unweighted form $\beta\wedge p^*\eta$ is top-degree and is
presented in $(b_i,z)$.  Corollary 4.9(2) puts their supports in
$S_i$ and $S_i\times\mathbf R$, respectively.

Multiplication by $\lambda_i$ or $p^*\lambda_i$ cannot enlarge
these supports.  Their paired supports are compact: the relevant
base support is compact inside $V_i$, and the parameter values lie
in a fixed compact interval, hence in a compact analytic annulus.
Finitely many affinoid neighborhoods of the base support give
compact total paired support.  This compactness justifies the
integrals; it is not used to exclude any tropicalized base boundary.

The restriction $\lambda_i|_{S_i}$ need only be PWS.  Remark 4.3(1)
makes its multiplication coefficientwise, and Proposition 2.17
gives the coefficientwise projection formula
\[
 (p_{S_i})_*\!\left(p_{S_i}^*(\lambda_i|_{S_i})\,T\right)
                  =(\lambda_i|_{S_i})(p_{S_i})_*T
\]
for the relevant polyhedral current $T$.
Use Definition 4.10(1) to evaluate the analytic top integrals on
these containing skeletons, followed by (2.17) and the integrable
current pushforward formula in Definition 4.10(4).  With
\eqref{r11p:product-test-skeletal}, this gives
\[
\begin{split}
 \int_{U\times\mathbf G_m}p^*\lambda_i\,\beta\wedge p^*\eta
 &=\int_{S_i\times\mathbf R}
      p_{S_i}^*(\lambda_i|_{S_i})
          (\beta\wedge p^*\eta)|_{S_i\times\mathbf R}\\
 &=\int_{S_i}(\lambda_i|_{S_i})(\theta\wedge\eta)|_{S_i}\\
 &=\int_U\lambda_i\,\theta\wedge\eta.
\end{split}
\]
Definition 4.10(2) permits summing this identity over the ordinary
partition.  It proves \eqref{r11p:product-pairing} for every
compact actual test.  The partition is not required to have a
single Euclidean presentation in the preassigned tuple $b_i$.

Finally, Euclidean proper-on-support pushforward commutes with
$d',d''$, and differentiation retains the compact parameter support.
Applying this to \eqref{r11p:product-definition}, whose actual
independence is already proved, gives the asserted differential
compatibility.  Equivalently it follows from the linear-map
current identity of Definition 2.29(3) in
\eqref{r11p:actual-product-skeletal}.
\end{proof}

\begin{lemma}[One-parameter evaluation of a nonlinear PL scalar]
\label{r11p:one-parameter}
Let $u$ be a bounded PL scalar on an ordinary open $V$.  Choose
$\chi\in C_c^\infty(\mathbf R)$ equal to one near its range.  On
$V\times(\mathbf G_m)^{\rm an}$ put
\[
                         H=\max(u-t,0).
\]
For every actual form $C$ on $V$ with a common moment presentation,
\[
                  I_p\bigl(\chi(t)t\,DH\wedge p^*C\bigr)=uC.
\tag{MT8}\label{r11p:graph-evaluation}
\]
Only one parameter occurs, independently of the number of source
moments needed to present $u$.
\end{lemma}

\begin{proof}
Write $u=\varphi(\ell)$ for a finite Euclidean PL scalar.  Work first
on its Euclidean moment domain with $H(x,t)=\max(\varphi(x)-t,0)$.
On each cell where $\varphi$ is affine, the graph face $t=\varphi(x)$
in $DH$ has projection weight one.  The coefficient of $t$ in the
defining affine difference is $-1$, so projection of its calibrated
divisor measure is the standard base measure.  Weighting by $\chi(t)t$
gives $\chi(\varphi(x))\varphi(x)$.

There are also vertical faces above walls of $\varphi$, in the region
$t<\varphi(x)$.  It would be false to identify $DH$ with the graph
alone.  On every such vertical face, however, projection drops the
dimension of its pure carrier; its pure Euclidean pushforward is
zero, exactly as in Mihatsch's Example 2.5.  Thus
\[
               \pi_*\bigl(\chi(t)t\,D\max(\varphi-t,0)\bigr)
                                  =\chi(\varphi)\varphi.
\]
The equality on affine cells extends across walls as an equality of
actual scalar functions; the right side is continuous and the
Euclidean pushforward is a delta-form of bidegree zero.

The affine projection formula now permits an arbitrary delta-form
$C$:
\[
 \pi_*\bigl(\chi(t)t\,DH\wedge\pi^*C\bigr)
       =\pi_*(\chi(t)t\,DH)\wedge C.
\]
Apply Lemma~\ref{r11p:product-integration} after moment refinement.
Only after actual moment pullback do we use the cutoff choice:
$\chi(\varphi(\ell))\varphi(\ell)=\chi(u)u=u$.
The ambient extension $\varphi$ need not be bounded, so
$\chi(\varphi)\varphi=\varphi$ is not asserted on its entire
Euclidean domain.
This argument has not pulled $C$ through the nonlinear PL graph;
only the affine projection pullback and projection formula are used.
\end{proof}

\subsection{The minimal finite trace theorem}

\begin{theorem}[Finite trace on the scalar-curvature module]
\label{r11p:minimal-trace-theorem}
Let $q:Y\to X$ be finite dominant between normal integral separated
varieties over $K$.  Every $\omega\in\mathscr M^k(Y^{\rm an})$ has a
unique actual trace on $X^{\rm an}$ satisfying
\[
                 \int_X\operatorname{Tr}^{\mathscr M}_q(\omega)
                                      \wedge\eta
                          =\int_Y\omega\wedge q^*\eta
\tag{MT9}\label{r11p:minimal-pairing}
\]
for all compact complementary actual tests.  It is local on the
ordinary topology, respects actual relations, and retains the full
native generic degree.  For a locally presented generator
$\omega=u_0Du_1\cdots Du_k$, a local formula is
\[
 \boxed{\displaystyle
 \operatorname{Tr}^{\mathscr M}_q(\omega)
 = I_{p_X}\!\left(
     \chi(t)t\,
     \operatorname{Tr}^{\rm curv}_{q\times\operatorname{id}}
       \left(D\max(u_0-t,0)\wedge\prod_{i=1}^kDu_i\right)
                         \right).}
\tag{MT10}\label{r11p:trace-formula}
\]
\end{theorem}

\begin{proof}
Shrink around an entire finite source fiber so that finitely many
generator presentations are combined as in
Lemma~\ref{r11p:pl-descent}, and $u_0$ is bounded there.  Choose
$\chi=1$ near that whole range.  The curvature trace inside
\eqref{r11p:trace-formula} exists by Lemma~\ref{r11p:pure-trace}.
Its parameter-uniform assertion supplies a product-moment presentation,
so multiplication by $\chi(t)t$ and application of $I_{p_X}$ are
legitimate.

Pair with an arbitrary compact actual test $\eta$.  First
Lemma~\ref{r11p:product-integration}, then
\eqref{r11p:pure-pairing} with the actual downstairs multiplier
$\chi(t)t\,p_X^*\eta$, and finally
Lemma~\ref{r11p:one-parameter}, give
\[
\begin{split}
 \int_X\operatorname{Tr}^{\mathscr M}_q(\omega)\eta
 &=\int_{X\times\mathbf G_m}\chi(t)t\,
     \operatorname{Tr}^{\rm curv}_{q\times\operatorname{id}}
        (DH\prod_iDu_i)\,p_X^*\eta\\
 &=\int_{Y\times\mathbf G_m}\chi(t)t\,DH\prod_iDu_i\,
                                          p_Y^*q^*\eta\\
 &=\int_Yu_0\prod_iDu_i\,q^*\eta.
\end{split}
\]
All these integrals have compact paired support.  The identity is
proved for all actual tests, not merely for model or original-smooth
tests.  Add generators.  If an expression represents zero as an
actual source form, its trace has zero pairing with every test, and
is zero by Proposition 4.11.  The same argument removes dependence
on the cutoff, source moments, quotient moments, Galois closure, and
all local choices.  It also glues the local answers in the ordinary
topology and proves uniqueness.
\end{proof}

\paragraph*{Limits of this theorem.}
It is enough for a PL scalar times a model Chern polynomial, since
model line curvatures locally have PL potentials.  It does not assert
a trace for an arbitrary source delta-form, and it does not assert
an analytic proper pushforward for a positive-dimensional proper map.
Normality of the target is retained throughout.  The Galois quotient
and native calibration arguments do not replace the original field
by an algebraic closure.  The theorem is independently derived above;
its use should be labelled as this internal minimal trace, rather than
as an externally published finite-trace result.

\paragraph*{Primary references.}
\begin{itemize}
\item \url{https://arxiv.org/html/2112.10018v4}
\item \url{https://arxiv.org/html/2107.12067v3}
\end{itemize}


\section{Finite orbit coordinates and a uniform parameter}

\begin{lemma}[An elementary polyhedral estimate]\label{h22:hoffman}
Let $C$ be a closed convex polyhedral cone in a finite-dimensional
Euclidean space and let $L$ be linear. There is a constant $c$ such that
\[
 \dist(z,C\cap\ker L)\leq c\|Lz\|\qquad(z\in C).
\]
\end{lemma}
\begin{proof}
Write $C=\{z:Az\geq0\}$, allowing pairs of opposite inequalities to
encode linear equalities. Let $y$ be the metric projection of $z$ onto
$Q=C\cap\ker L$, put $u=z-y$, and let $I$ index the inequalities active
at $y$. Thus $A_Iu\geq0$. The elementary normal-cone formula for a
polyhedron gives
\[
 u\in N_Q(y)=\{-A_I^{\mathsf T}\lambda+L^{\mathsf T}\mu:
                                  \lambda\geq0\}.
\]
This formula is the finite-dimensional polar-cone alternative for the
linear inequalities $A_Iv\geq0$, $Lv=0$: a vector is in the normal cone
exactly when its scalar product with each such tangent vector is
nonpositive. Define the closed polyhedral cone
\[
 E_I=\{u:A_Iu\geq0\}\cap
     \{-A_I^{\mathsf T}\lambda+L^{\mathsf T}\mu:\lambda\geq0\}.
\]
If $u\in E_I$ and $Lu=0$, its displayed representation gives
\[
 \|u\|^2=-\lambda^{\mathsf T}A_Iu+\mu^{\mathsf T}Lu\leq0,
\]
so $u=0$. Compactness of $E_I\cap\{\|u\|=1\}$ therefore gives
$\|u\|\leq c_I\|Lu\|$ on $E_I$; an empty intersection causes no
problem. There are finitely many possible active sets. Taking their
maximum and using $Ly=0$ proves the assertion.
\end{proof}

\begin{lemma}[PL factorization through finite orbit coordinates]
\label{h22:orbit-pl}
Let a finite group $G$, of order $d$, act on $\R^s$ by coordinate
permutations, where $s\geq1$. Put $N=d(s-1)+1$, choose distinct integers
$a_1,\ldots,a_N$, and set
\[
 m_j=(1,a_j,\ldots,a_j^{s-1}).
\]
Let $J:\R^s\longrightarrow\R^{Nd}$ record, for each $j$, the ordered
multiset $\{\langle m_j,gx\rangle:g\in G\}$. Then:
\begin{enumerate}
\item $J(x)=J(y)$ if and only if $x=gy$ for some $g\in G$.
\item There is a constant $C$ such that
\[
 \min_{g\in G}\|x-gy\|\leq C\|J(x)-J(y)\|.
\]
\item For every $G$-invariant finite PL function
$\varphi:\R^s\times\R\to\R$, where $G$ fixes the last coordinate,
there is a finite PL function $\Phi:\R^{Nd}\times\R\to\R$ with
\[
                   \varphi(x,t)=\Phi(J(x),t).
\]
The same assertion holds for any finite collection of functions with
one common map $J$.
\end{enumerate}
\end{lemma}
\begin{proof}
If $J(x)=J(y)$, for each $j$ choose $g_j\in G$ with
$\langle m_j,x-g_jy\rangle=0$. Some $g$ occurs at least $s$ times.
The polynomial with coefficients $x-gy$ has degree at most $s-1$ and
has these $s$ distinct roots. It is zero, so $x=gy$. The converse is
immediate.

The ordering hyperplanes give a finite fan of closed convex polyhedral
cones $C_\alpha$ on each of which $J=A_\alpha$ is linear. For a pair
of cones apply Lemma~\ref{h22:hoffman} to
\[
 C_\alpha\times C_\beta,
 \qquad L(x,y)=A_\alpha x-A_\beta y.
\]
If $(x_0,y_0)$ is in its kernel, part~(1) gives $x_0=gy_0$ for
some $g\in G$. Since the action is orthogonal,
\[
 \min_g\|x-gy\|\leq
 \|x-x_0\|+\|y-y_0\|
 \leq\sqrt2\|(x,y)-(x_0,y_0)\|.
\]
The finitely many cone-pair estimates prove part~(2).

Write $A=J(\R^s)$ and define $\psi(J(x),t)=\varphi(x,t)$ on
$A\times\R$. This is well defined. A finite PL function on a Euclidean
space is globally Lipschitz, by the finite bound on its affine gradients
and subdivision of a straight segment. Part~(2), applied after choosing
the best orbit representative, shows that $\psi$ is globally Lipschitz
for the induced Euclidean metric on $A\times\R$.

The map $A_\alpha$ is injective on the span of $C_\alpha$: otherwise a
nonzero kernel vector would give infinitely many points with the same
$J$-value near a relative interior point, contradicting part~(1).
Refine $C_\alpha\times\R$ by the affine cells of $\varphi$. Their
images under $J\times\Id$ form a finite closed polyhedral cover of
$A\times\R$ on which $\psi$ is affine. Values agree on overlaps.

Here is an explicit ambient extension. Choose $L_0$ at least a
Lipschitz constant for $\psi$ with respect to the $\ell^1$ norm and put
\[
 \Phi(w)=\inf_{a\in A\times\R}
              \bigl(\psi(a)+L_0\|w-a\|_1\bigr).
\tag{TD1}\label{h22:mcshane}
\]
The Lipschitz inequality shows that the infimum is finite and that
$\Phi=\psi$ on $A\times\R$. It also shows that $\Phi$ is Lipschitz.
On each of the finitely many polyhedra where $\psi$ is affine, \eqref{h22:mcshane} is
the value of a finite linear program: introduce variables $r_i$ with
$r_i\geq w_i-a_i$ and $r_i\geq a_i-w_i$ and minimize the affine
function $\psi(a)+L_0\sum_i r_i$. Its epigraph is the projection of a
polyhedron, hence a polyhedron. The value is everywhere finite and is
therefore a finite convex PL function. Taking the minimum over the
finite cover again gives a finite continuous PL function. This proves
part~(3), including its whole-parameter-line assertion. The same $J$
serves every member of a finite collection; their Lipschitz constants
can be bounded by one common constant if desired.
\end{proof}

\begin{proposition}[Ordinary scalar descent with fixed parameter moments]
\label{h22:scalar-descent}
Let $K$ be complete, nontrivially valued and of characteristic zero.
Let $q:Y\to X$ be a finite dominant morphism of normal integral
separated algebraic $K$-varieties whose function-field extension is
Galois with group $G$. A $G$-invariant PL scalar on $Y^{\an}$ descends
locally in the ordinary topology to a PL scalar on $X^{\an}$.

For the parameter version, fix $x\in X^{\an}$ and an ordinary
neighborhood $U$ of $x$. Suppose finitely many invariant scalar
functions on $q^{-1}U\times\mathbf G_m^{\an}$ admit presentations
\[
 F_i(y,z)=\varphi_i(\ell(y),t(z)),\qquad
 \ell=(v(f_1),\ldots,v(f_s)),\quad t=v(z),
\]
where $f_j$ are a fixed tuple of analytic units on $q^{-1}U$ and
$\varphi_i$ are finite PL functions on the whole Euclidean moment
space. After shrinking $U$ around $x$, there are a fixed tuple of base
analytic units $b$ and finite PL functions $\Psi_i$ such that
\[
                  F_i=q^*\Psi_i(v(b),t).
\]
The base tuple and shrinking are independent of every $t\in\R$.
\end{proposition}
\begin{proof}
For the nonparameter assertion, choose disjoint relatively compact
ordinary neighborhoods of the finitely many points of $q^{-1}(x)$,
small enough for moment presentations of the prescribed scalars.
Properness of $q^{\an}$ gives an ordinary neighborhood $U$ of $x$
whose full inverse image lies in their union. The source pieces are
clopen in $q^{-1}U$. Paste analytic units across these pieces, inserting
constant units on the other pieces. Use one additional pasted constant
unit with different valuations to distinguish the pieces. Nontriviality
of the valuation supplies arbitrarily many such constants. By the
initial shrinking the source moment images lie in finitely many
disjoint compact boxes. A finite triangulation of a surrounding compact
box, subordinate to these boxes and to the scalar affine cells,
extends each of their finite PL functions. Extending beyond the
surrounding box by coordinatewise clipping gives one finite PL
presentation on the full source moment space.

In the parameter case the stipulated common presentation is the
starting point; no extension from arbitrary unbounded disjoint
cylinders is needed. Enlarge the source tuple by its $G$-orbit.
The group now acts by coordinate permutations. Average each ambient
presentation over $G$. Invariance on the source shows that averaging
does not change the represented scalar function, while it produces
an invariant finite PL function on the whole Euclidean space.
Apply Lemma~\ref{h22:orbit-pl}, with or without the parameter.

It remains to express the coordinates of $J(\ell)$ in base moments.
For the integer direction $m_j$ let $f^{m_j}$ be the corresponding
analytic unit. Form the polynomial
\[
                    P_j(T)=\prod_{g\in G}(T-g(f^{m_j})).
\]
Its coefficients are invariant analytic functions and descend to $U$.
One can check this last assertion directly on an affinoid domain
inside an affine algebraic chart $\operatorname{Spec}R\subset X$.
If $\operatorname{Spec}S=q^{-1}(\operatorname{Spec}R)$, normality
gives $S^G=R$: an invariant element lies in $\operatorname{Frac}(R)$
and is integral over $R$. For an affinoid domain with algebra $A$, its
finite inverse image has algebra $S\otimes_R A$. The averaging
idempotent on $S$ splits $S=R\oplus\ker(e)$. Tensoring the splitting
with $A$ shows $(S\otimes_R A)^G=A$. This also proves descent on
ordinary overlaps.

At $x$, retain every nonzero coefficient of each $P_j$, including its
constant coefficient, which is a unit. Shrink $U$ so that the retained
coefficients are analytic units. If the constant coefficient is
$a_0$, arrange $v(a_0)\leq B$. A coefficient vanishing at $x$ has,
after further shrinking, valuation greater than $\max(B,0)+1$.
Its point in the Newton diagram lies above the chord from
$(0,v(a_0))$ to $(d,0)$ and cannot affect the lower hull. Thus all
discarded coefficients may be omitted uniformly on $U$.

The ordered root valuations of a fixed-degree polynomial are finite
PL functions of the valuations of its retained coefficients: on each
region specifying its lower-hull vertices, they are the negatives of
the intervening slopes, with the appropriate multiplicities. At every
point over the base these roots are the orbit values used to define
$J$. Consequently $J(\ell)$ is a finite PL function of a single finite
tuple $v(b)$ of base-unit valuations. These coefficients involve only
the original source units and are independent of $t$. Composing this
scalar PL expression with the extensions from
Lemma~\ref{h22:orbit-pl} proves the assertion. Derivatives are taken
only after this scalar construction; no nonlinear pullback of a
general delta-form is used.
\end{proof}

\begin{remark}[Why the parameter hypothesis matters]\label{h22:parameter-scope}
There is no assertion here that arbitrary finite PL functions given
on disjoint unbounded polyhedra always admit a finite PL ambient
extension. That assertion is false: on the two parallel lines
$\{0\}\times\R$ and $\{1\}\times\R$, prescribe respectively the
values $t$ and $0$. A finite PL function on $\R^2$ is globally
Lipschitz, so cannot have these restrictions. The parameter statement
above uses its specified common finite PL presentation. In the minimal
trace application, first combine the base potentials $u_0,\ldots,u_k$
on compact source moment boxes, and only then form
$\max(u_0-t,0)$. This produces exactly the required common finite
jointly PL presentation for the whole parameter line.
\end{remark}

\section{A compact-parameter alternative sufficient for the trace}

\begin{proposition}[Uniform moments only on the necessary annulus]
\label{h22:compact-parameter}
In the formula for the minimal finite trace of
$u_0Du_1\cdots Du_k$, the whole-line parameter assertion of
Proposition~\ref{h22:scalar-descent} may be replaced by its assertion
on one fixed compact interval $I$ containing the support of the cutoff
in its interior. The resulting trace has exactly the same domain,
all-actual pairing property, native degree and independence of
presentations as the original formula.
\end{proposition}
\begin{proof}
Shrink the base around the finite fiber so that the base potentials
have common bounded moment boxes. Choose $\chi\in C_c^\infty(\R)$
equal to one near the range of $u_0$, and choose a compact interval
$I$ with $\Supp\chi\subset I^\circ$. All source boxes times $I$
are compact. Finite PL functions on their finite polyhedral union
can be extended by a finite triangulation of a larger compact box.
The same is true for the compact image of the orbit-coordinate map:
subdivide by its injective linearity cells and by the given scalar
affine cells; invariant values agree on overlapping images. Thus the
scalar-descent proof needs only compact finite polyhedral extension,
followed by the same Newton-polynomial coefficient construction.

Sorting and polarization give a fixed product-moment presentation
$\Theta$ for the pure-curvature trace on
$U\times\{t\in I^\circ\}$, after one common base shrinking.
Its ambient Euclidean coefficient may be extended outside that
interval by the finite PL potential extensions just constructed.
Multiplication by $\chi(t)t$ kills every possible discrepancy near
and outside the ends of $I$. Hence $\chi(t)t\Theta$ has a genuine
global product-moment presentation with compact parameter support
on $U\times\mathbf G_m^{\an}$. It is the zero extension of its
restriction on the annulus. No extension of an arbitrary current or
arbitrary delta-form is required.

The product-integration argument applies to this compactly supported
coefficient, and its pairing with any compact actual base test is
computed entirely on that annulus. The one-parameter graph evaluation
likewise sees only $\Supp\chi$. It therefore returns
$u_0Du_1\cdots Du_k$ with exactly the same pairing formula as in
Theorem~\ref{r11p:minimal-trace-theorem}. Detection by ordinary smooth tests, a subset of the actual
tests already allowed, removes all choices and glues the resulting
actual forms. Neither the final trace theorem nor any of its tests
has been restricted.
\end{proof}

\section{The product skeleton and its topology}

\begin{proposition}[Coordinate Gauss product]
\label{h22:gauss-product}
Let $V$ be a boundaryless ordinary analytic open of pure dimension
$n$ in an algebraic $K$-analytic space, and let $b$ be a finite tuple
of analytic units on $V$. Put $S=\Sigma(V,b)$. The relative Gauss
section identifies
\[
 S\times\R\ \simeq\
 \Sigma(V\times\mathbf G_m^{\an},(b,z)),\qquad
              (y,t)\longmapsto\nu_t(y),
\tag{TD2}\label{h22:gauss}
\]
as weighted pwl spaces with their specified linear function sheaves.
The weight is the base weight times the standard weight on $\R$.
The identification respects ordinary restriction and every additional
base moment tuple.
\end{proposition}
\begin{proof}
The skeleton criterion is the graded-transcendence-degree criterion
of \cite[\S3.1]{MihLT}. A point above $y$ has degree $n+1$ in the
moments $(b,z)$ if and only if $b$ has degree $n$ at $y$ and the
graded reduction of $z$ is transcendental over
$\widetilde{\mathcal H(y)}$. For this equivalence, once the base
degree is $n$, the full graded residue field of $y$ is algebraic over
the field generated by $b$. The last condition means precisely that
every Laurent polynomial has the Gauss norm
\[
 \left|\sum_i a_i z^i\right|=\max_i |a_i(y)|e^{-it}.
\]
It thus specifies the unique point $\nu_t(y)$. This proves the
bijection in \eqref{h22:gauss}, without extending $K$ or restricting the radius.

Here are the topology details. On an affinoid base domain and a
compact annulus in the parameter, analytic functions have uniformly
convergent Laurent expansions. Their values on the Gauss section
are limits of the finite maxima displayed above, uniformly on that
domain times the parameter interval. They are continuous in $(y,t)$.
The Gauss section is therefore continuous locally, hence globally.
Its inverse is the continuous map $(p,v(z))$ on its image.
Equivalently, its image over a compact base set and a compact interval
is closed inside the corresponding compact relative annulus, since
the displayed Gauss equalities are closed conditions and determine
the image uniquely. This verifies, when needed, the properness claim
used in the product-integration argument.

The inverse map is pwl for the skeleton atlases: a pwl chart on a
base skeleton is locally described by valuations of analytic units
on an affinoid domain; their pullbacks, together with $v(z)$, give
the same coordinate functions upstairs. Analytic-domain restrictions
and their inverse images restrict the corresponding pwl charts.
After taking finite polyhedral charts, this continuous bijective pwl
map is a pwl isomorphism. Indeed subdivide until each piece is
affine. Injectivity makes each linear part injective on the span of
its piece, so its inverse is affine on the image. Properness on
compact charts and a common image subdivision make these inverse
pieces a locally finite pwl atlas. This is also the coordinate case
of the atlas argument following (3.31) in \cite{MihLT}.

On a top-dimensional base cell, compute its native weight by a
generic toric projection that is locally finite flat, as in the
definition of the skeleton weight. Its product with the identity
on the parameter torus is finite flat of the same degree. Hence the
weight on the product cell is the base weight times the standard
parameter weight. All residue and field-degree factors are retained.
The moment tuple is literally $(b,z)$, so the specified sheaf of
linear functions is generated by the base linear functions and $t$.
Adding base moments uses the same relative Gauss section and the
same degrees; restriction does likewise. This proves all the claimed
compatibilities.
\end{proof}

\section{How the two details enter the full-test trace argument}

Let $q:Y\to X$ be finite dominant between the normal integral
varieties of Section~\ref{r11p:minimal-trace}. The following records
how scalar descent and the coordinate Gauss product enter the
full-test trace formula.

For PL scalars $u,v$, set $h=u-v$, $h_+=\max(h,0)$ and
$h_-=\min(h,0)$. Coefficientwise multiplication on every witnessing
skeleton gives
\[
 h_+Dh_-=h_-Dh_+=0.
\]
The actual Leibniz rule and separate closedness then give
$Dh_+\wedge Dh_-=0$. Expanding proves the max/min curvature power
identity. A finite sorting network and polarization descend the sum
of conjugate pure curvature products. Proposition~\ref{h22:scalar-descent},
or Proposition~\ref{h22:compact-parameter} in the cutoff application,
provides its required scalar presentations. A normal Galois closure
and division by its full generic degree handle the nongalois case.

The native top-degree identity
\[
                  \int_Yq^*\tau=(\deg q)\int_X\tau
\]
follows from generic flatness, support avoidance of the nonflat
algebraic locus, and the native skeletal finite-flat degree identity.
It proves the pure-curvature pairing identity and faithfulness of
finite pullback; it does not assert pushforward of arbitrary source
forms.

For a product-moment form $\beta=t_{(c,z)}^*\Gamma$ with compact
parameter support, let $\pi$ denote the Euclidean coordinate
projection and define
\[
                         I_p(\beta)=t_c^*(\pi_*\Gamma).
\]
Proposition~\ref{h22:gauss-product} identifies the relevant product
skeletons. On each finite-fiber tropical chart $g$ and each compact
polyhedral set $H$, keep the excluded set
$E_H=g(\partial H)$. On the complement of $E_H$, the characterization
(2.40) of \cite[Proposition~2.38]{MihLT} and the affine projection
formula of \cite{MihTrop} identify the pushforward on the product
skeleton with $t_c^*(\pi_*\Gamma)$. Varying $H$ detects the whole
polyhedral current by \cite[Lemma~2.37]{MihLT}. The boundary exclusion
must not be dropped.

For an arbitrary compact actual test $\eta$, concatenate its moment
tuple with $c$. The same identity then applies to the Euclidean
coefficient of $\beta\wedge p^*\eta$. Original smooth partitions
give compact local paired supports, and intrinsic skeletal integration
gives
\[
                   \int_X I_p(\beta)\wedge\eta
                    =\int_{X\times\mathbf G_m}\beta\wedge p^*\eta.
\]
In particular the test is not required to be an external product,
smooth, or a polynomial in model curvatures.

Finally $D\max(\varphi(x)-t,0)$ has graph faces, with projection
weight one, and vertical faces over the bends of $\varphi$. Projection
annihilates the latter pure carriers because it drops their dimension.
This yields
\[
 \pi_*\bigl(\chi(t)tD\max(\varphi(x)-t,0)\bigr)
                         =\chi(\varphi(x))\varphi(x).
\]
The affine projection formula permits multiplication by an arbitrary
actual coefficient after moment refinement. Thus, for
$\chi=1$ near the range of $u_0$,
\[
 \Tr_q^{\mathscr M}(u_0Du_1\cdots Du_k)
  =I_{p_X}\left(\chi(t)t\,
    \Tr_{q\times\Id}^{\mathrm{curv}}
       \left(D\max(u_0-t,0)\prod_{i=1}^kDu_i\right)\right)
\]
satisfies the full complementary-actual-test pairing. Detection by
ordinary smooth tests proves that actual source relations have zero
trace, and removes the cutoff, Galois-closure, moment and local choices.
This is exactly the minimal finite trace required for the supported
perfect-complex construction.



\section{Supported complexes, prescribed cycles, and global primitives}
\label{r11p:supported-complex}

\paragraph*{Construction and dependencies.}
The supported-complex construction uses the actual model Chern forms of
Theorem~\ref{h22c:chern} and the one-parameter minimal trace of
Theorem~\ref{r11p:minimal-trace-theorem}. A smooth finite multisection
is the algebraic ingredient; its role and proof are stated below.

Let $K$ be complete discretely valued of characteristic zero, $R=K^\circ$,
and let $\mathcal X/R$ be integral, flat and projective, with smooth
integral projective generic fiber $X$ of dimension $n$.
The special fiber may be singular or nonreduced.  Let
\[
       C=(0\longrightarrow C_N\longrightarrow\cdots
                             \longrightarrow C_0\longrightarrow0)
\]
be a bounded complex of model vector bundles, exact on
$\mathcal U=\mathcal X\setminus\mathcal Z$, with generic open
$U=X\setminus Z$.  Assume $Z$ is pure of codimension $r$, $1\leq r\leq n$.
The same construction is linear in finite rational combinations of
such complexes.  Define
\[
 \Omega_C=\sum_i(-1)^i\operatorname{ch}_r(\overline C_i)
                              \in\mathcal B^{0,0,r}(X^{\rm an}),
 \qquad
 z_C=\operatorname{ch}^{Z}_r(C|_X)\cap[X]
                              \in\operatorname{CH}_{n-r}(Z)_{\mathbb Q}.
\tag{PC1}\label{r11p:omega-and-cycle}
\]
The first expression is an actual form by Theorem~\ref{h22c:chern}.
The second is a cycle with coefficients, because
$\operatorname{CH}_{n-r}(Z)$ is the free group on the top-dimensional
components of $Z$.  It is not merely an ambient Chow class.

\subsection{The algebraic finite multisection actually needed}

\begin{lemma}[A smooth finite flat multisection]
\label{r11p:multisection}
Let $F\to Y$ be a smooth projective surjection of smooth projective
integral varieties over the infinite field $K$, of pure relative
dimension $d$.  There is a smooth projective integral closed subscheme
$T_0\subset F$ whose map to $Y$ is finite flat surjective.
It may be taken as a connected component of a smooth complete intersection.
If the source is geometrically integral and $\dim Y\geq1$, $T_0$
can be chosen geometrically integral.
\end{lemma}

\begin{proof}
Choose a very ample $H$ on $F$ and use $M=H^m$ with
$m\geq\max(\dim Y,1)$.  For any positive-dimensional integral component
$A$ of any intermediate fiber complete intersection, two sections
$a_0,a_1$ of $H$ have nonconstant ratio on $A$.  Consequently
\[
                    a_0^m,a_0^{m-1}a_1,\ldots,a_1^m
\]
are linearly independent on $A$.  In the parameter space
$V=H^0(F,M)$, the sections vanishing identically on $A$ therefore
have codimension at least $m+1$.

Fix a geometric fiber, which is pure Cohen--Macaulay of dimension $d$.
Induct on $j\leq d$ to see that the bad $j$-tuples, whose common
zero locus has dimension greater than $d-j$, have codimension at least
$m+1$ in $V^j$.  Above a good prefix, a bad next section vanishes
on one of the finitely many components just considered; above a bad
prefix the inductive bound applies.  Upper semicontinuity of fiber
dimension makes the universal bad locus closed.  Adding $\dim Y$
still leaves its dimension strictly smaller than $\dim V^d$.
Properness of $Y$ makes its projection to $V^d$ closed.  Thus a
nonempty open of $d$-tuples has zero-dimensional fibers everywhere.
The fibers are nonempty, because $d$ ample divisors meet every
nonempty pure $d$-dimensional projective fiber.

Smooth Bertini supplies a second nonempty open on which the total
intersection is smooth of dimension $\dim Y$.  Geometric irreducibility
Bertini supplies the stated refinement when relevant.  Their intersection
has a $K$-rational point since $K$ is infinite.  The resulting proper
quasi-finite morphism is finite and surjective.  Its source is
Cohen--Macaulay and its target regular of the same dimension, so
miracle flatness applies.  If necessary choose a connected component;
it is integral, finite flat, and still surjective.
\end{proof}

The bound $m+1>\dim Y$, not just Bertini, is what rules out a
positive-dimensional fiber anywhere on the base.  The primary
ingredients are Stacks Tags 0D4I, 0FD6, 0G4C and 00R4.

\begin{lemma}[Relative diagonalization with the original open fixed]
\label{r11p:relative-lines}
There is a projective morphism of integral flat models
$f:\mathcal T\to\mathcal X$ with smooth projective generic fiber $T$,
such that $f:T\to X$ is generically finite of degree $d$ and finite
flat over the \emph{entire} prescribed open $U$.  There are model
line maps $s_\alpha:L_\alpha\to M_\alpha$, isomorphisms above
$\mathcal U$, and signs $\epsilon_\alpha$, satisfying
\[
                 [f^*C]
                   =\sum_\alpha\epsilon_\alpha[L_\alpha\to M_\alpha]
\tag{PC2}\label{r11p:supported-k-equality}
\]
in supported perfect $K_0$.  Their generic divisors
$D_\alpha=\operatorname{div}(s_\alpha|_T)$ are supported in $f^{-1}Z$.
\end{lemma}

\begin{proof}
On $\mathcal U$, exactness of a bounded vector-bundle complex is locally
split.  Hence $S_i^\circ=\ker d_i=\operatorname{im}d_{i+1}$ are
subbundles, with $S_0^\circ=C_0$ and $S_N^\circ=0$.  Close the graph
of this section in the product of the corresponding Grassmann bundles.
This gives a projective birational integral model
$\mathcal Y\to\mathcal X$, an isomorphism over $\mathcal U$, with
tautological subbundles $S_i\subset C_i$ and quotient bundles $Q_i$.
Every asserted vanishing of a bundle map that holds on the dense open
holds on this integral model.  The differential therefore factors as
\[
                            a_i:Q_i\longrightarrow S_{i-1},
\]
an equal-rank map that is an isomorphism over $\mathcal U$.

There is an actual filtration by subcomplexes: its $p$th term has
$C_i$ for $i>p$, $S_p$ in degree $p$, and zero below $p$.
The successive two-term quotients are
$Q_{p+1}\to S_p$ in degrees $p+1,p$.  They are exact on the given
open, and yield
\[
                     [C]=\sum_i(-1)^{i-1}[Q_i\to S_{i-1}]
\]
in supported $K_0$, rather than only in ordinary $K_0$.

Resolve the generic $\mathcal Y_K$ in characteristic zero, without
changing its smooth open $U$, and extend the projective resolution by
graph closure over $R$.  This uses resolution over the field $K$ only,
not mixed-characteristic resolution of a model.  Call the resulting
smooth projective generic fiber $Y$.

Take the product of the full source flag bundles of \emph{all} $Q_i$
over $Y$.  Before extending the transported target flags, choose a
smooth finite flat multisection $T_0\to Y$ by
Lemma~\ref{r11p:multisection}.  Take its flat model closure.  Over
$\mathcal U$, the isomorphisms $a_i$ transport every source flag to a
target flag.  Close the graph of these target flags, and resolve the
resulting generic fiber without changing its smooth inverse image of
$U$.  Its final generic fiber $T$ is therefore still exactly $T_0$
over $U$.  In particular $T\to X$ is finite flat over all of $U$;
the successive modifications have introduced no new excluded image.

The maps $a_i$ preserve the two full flags.  Their associated graded
maps are model line maps $L_\alpha\to M_\alpha$, isomorphisms on the
original model open.  Termwise filtration of the two-term complexes
gives \eqref{r11p:supported-k-equality}.  A nonzero map of line bundles
on the integral model is injective and defines an effective Cartier
divisor.  Its generic divisor can only lie above $Z$.  Empty generic
divisors are allowed, for instance for a purely vertical line map.
\end{proof}

\paragraph*{Why ordinary flag sections alone do not replace this step.}
The model-Chern construction in Section~\ref{h22c:section} descends forms by local sections because the
upstairs Chern form satisfies its fiber-product cocycle identity.
A primitive on a flag bundle need not satisfy that identity: pulling
it to two sections only produces forms whose difference is $D$-closed.
The present argument does not infer Green descent from Chern descent.
Its use of a finite multisection is exactly where the minimal trace is
subsequently applied.

\subsection{The source cycle with all coefficients retained}

For degree-one variables $l,m$ put
\[
 B_{r-1}(l,m)=\frac1{r!}\sum_{a=0}^{r-1}m^{r-1-a}l^a,
 \qquad
                  (m-l)B_{r-1}(l,m)=\frac{m^r-l^r}{r!}.
\tag{PC3}\label{r11p:B-polynomial}
\]
For a nonzero line map $s:L\to M$ on $T$, with divisor $D$, the
all-degree Cartier localized character formula is
\[
 \operatorname{ch}^{D}_r([L\to M])\cap[T]
         =B_{r-1}(c_1L,c_1M)\cap[D].
\tag{PC4}\label{r11p:cartier-localized}
\]
Here $L$ has degree one and $M$ degree zero.  To check the full
polynomial, the complex resolves $M|_D$ and the normal bundle of
$D\hookrightarrow T$ is $(M\otimes L^{-1})|_D$.  Localized
Riemann--Roch for this regular immersion gives
\[
 e^m\,\frac{1-e^{-(m-l)}}{m-l}\cap[D]
                         =\frac{e^m-e^l}{m-l}\cap[D].
\]
Its homogeneous degree $r-1$ is exactly \eqref{r11p:cartier-localized}.
This is Fulton, \emph{Intersection Theory}, Corollary 18.1.2,
specialized to a Cartier immersion; the proof uses deformation to the
normal bundle.  A leading-term statement alone, such as Stacks 0FEH,
would not justify this all-degree polynomial.

Additivity on supported perfect $K_0$, compatibility with arbitrary
base change, and the proper bivariant projection formula now give
\[
 \sum_\alpha\epsilon_\alpha
    (f|_{D_\alpha})_*
       \left(B_{r-1}(c_1L_\alpha,c_1M_\alpha)\cap[D_\alpha]\right)
                              =d\,z_C
       \quad\hbox{in }\operatorname{CH}_{n-r}(Z)_{\mathbb Q}.
\tag{PC5}\label{r11p:source-equality}
\]
The underlying equality is
$f_*(\operatorname{ch}^{f^{-1}Z}(f^*C)\cap[T])
=\operatorname{ch}^Z(C)\cap f_*[T]$, and $f_*[T]=d[X]$.
These are supported identities; one has not first forgotten the
support and then attempted to recover coefficients.  Nilpotent
thickenings can be used to make every indicated factorization through
$Z$ scheme-theoretic.  Chow groups are unchanged by those thickenings.
Since the target Chow group is top-dimensional on $Z$,
\eqref{r11p:source-equality} is equality of actual cycle coefficients.
See Stacks 0FB3 (base change), 0FBE (additivity), and 0FET (supported
perfect $K_0$).

\subsection{A proper degree calculation, only against top tests}

\begin{lemma}[Chern degree against an arbitrary actual top test]
\label{r11p:proper-chern-degree}
Let $h:B\to A$ be proper between integral projective $K$-varieties,
$\dim A=m$, $\dim B=m+t$.  Let $P$ be a homogeneous degree-$t$
polynomial in first Chern forms of algebraic line bundles on $B$
with PWS metrics.  For every compact actual top form $\eta$ on $A$,
\[
              \int_BP\wedge h^*\eta
                =\int_{h_*(P_{\rm alg}\cap[B])}\eta.
\tag{PC6}\label{r11p:proper-top-pairing}
\]
The statement is linear in cycles and includes their generic lengths.
\end{lemma}

\begin{proof}
Write each algebraic line bundle as a quotient of two very ample ones.
Choose PWS metrics on those factors so that their quotient is the
given metric.  Expand $P$ by multilinearity.  For each resulting
monomial, choose a section of each very ample factor successively,
meeting the finitely many current integral cycle components properly.
The infinite field $K$ supplies such choices.  Work componentwise
when an intermediate divisor is reducible or nonreduced.

At every step apply scalar strong Poincare--Lelong, Mihatsch
Corollary 5.6, with the remaining closed curvature factors and the
pullback of $\eta$ as test.  The latter is separately closed because
$\eta$ is top-degree on $A$.  The Green correction is therefore zero.
One is left with the pullback of $\eta$ integrated on a cycle of
dimension $m$.

A component mapping to a lower-dimensional subset of $A$ contributes
zero, because the support of the actual top test avoids that subset.
For a dominating component, remove a proper closed subset of $A$
so that the component maps finite flat of its native generic degree.
The top test again avoids the removed subset.  The native finite
degree formula evaluates the integral with exactly that degree.
This is precisely algebraic proper pushforward of the terminal
intersection cycle, proving \eqref{r11p:proper-top-pairing}.
\end{proof}

This lemma is an all-actual-test scalar equality.  It does not define
the proper pushforward of an arbitrary actual form.  It applies to
each $D_\alpha\to Z_j$ in \eqref{r11p:source-equality}, with
$m=n-r$ and $t=r-1$.  Components whose image has smaller dimension
contribute zero by the same support argument.

\subsection{The exact-open Green construction}

\begin{theorem}[A strong Green pair for a supported model complex]
\label{r11p:complex-green}
Under the assumptions of \eqref{r11p:omega-and-cycle}, there is an
actual form
\[
                       G_C\in\mathcal B^{r-1,r-1}(U^{\rm an})
\]
such that, for every compact actual
$\eta\in\mathcal B_c^{n-r,n-r}(X^{\rm an})$,
\[
                 \int_UG_C\wedge D\eta+\int_{z_C}\eta
                                     =\int_X\Omega_C\wedge\eta.
\tag{PC7}\label{r11p:complex-green-equation}
\]
The domain is the prescribed $U=X\setminus Z$, with no auxiliary
closed subset removed.  If a tube is disjoint from the full model
support $\mathcal Z$, the construction can be chosen zero on that tube.
\end{theorem}

\begin{proof}
Use Lemma~\ref{r11p:relative-lines} and give the line bundles their
model metrics.  Write
\[
 l_\alpha=c_1(\overline L_\alpha),\qquad
 m_\alpha=c_1(\overline M_\alpha),\qquad
 g_\alpha=-\log\|s_\alpha\|.
\]
On $f^{-1}U$ each $g_\alpha$ is a PL scalar.  Put
\[
       G_T=\sum_\alpha\epsilon_\alpha
                    g_\alpha B_{r-1}(l_\alpha,m_\alpha),
       \qquad
       G_C=\frac1d\operatorname{Tr}^{\mathscr M}_{f|_{f^{-1}U}}(G_T).
\tag{PC8}\label{r11p:complex-green-formula}
\]
Every summand lies in the scalar-curvature module of
Theorem~\ref{r11p:minimal-trace-theorem}.  Source and target of the
finite restriction are smooth, hence normal; all its assumptions hold.
The resulting form is actual on the full ordinary open $U$.

Strict model Chern additivity and functoriality from Theorem~\ref{h22c:chern} convert
\eqref{r11p:supported-k-equality}, after forgetting support, into
the actual identity on $T$
\[
 \sum_\alpha\epsilon_\alpha
             (m_\alpha-l_\alpha)B_{r-1}(l_\alpha,m_\alpha)
                                      =f^*\Omega_C.
\tag{PC9}\label{r11p:actual-curvature-equality}
\]
No model-to-cycle Bott--Chern comparison is being used here.

Apply scalar strong Poincare--Lelong to each $s_\alpha$ with test
$B_{r-1}(l_\alpha,m_\alpha)f^*\eta$.  These tests are compact because
$f$ is proper.  Curvature factors are separately closed.  Summation
and \eqref{r11p:actual-curvature-equality} yield
\[
 \int_{T\setminus f^{-1}Z}G_Tf^*D\eta
 +\sum_\alpha\epsilon_\alpha
       \int_{[D_\alpha]}B_{r-1}(l_\alpha,m_\alpha)f^*\eta
                         =\int_Tf^*(\Omega_C\eta).
\tag{PC10}\label{r11p:upstairs-complex-PL}
\]
For the first term, extending each off-divisor integral to the
displayed common open changes nothing: $D\eta$ has bidegree
$(n-r+1,n-r+1)$, so its closed support avoids $Z$ by
Mihatsch Corollary 4.9.  Properness puts its pulled-back support
inside the finite locus $f^{-1}U$.  The minimal trace pairing identifies
this first term with $d\int_UG_CD\eta$.

The form $\Omega_C\eta$ is top-degree on $X$, so its support also
avoids $Z$.  On that support $f$ is finite flat of degree $d$;
the last term is $d\int_X\Omega_C\eta$ by the native degree formula.
The middle term equals $d\int_{z_C}\eta$ by
Lemma~\ref{r11p:proper-chern-degree} and
\eqref{r11p:source-equality}.  Divide by $d$.

All tests defining the off-$Z$ integration functional of $G_C$
have bidegree $(n-r+1,n-r+1)$ and therefore support away from $Z$.
Thus its canonical extension and the all-actual meaning of
\eqref{r11p:complex-green-equation} require no integrability
assumption near $Z$ beyond those already encoded by support avoidance.

Finally, above the model open $\mathcal U$ each $s_\alpha$ is a model
isomorphism of line bundles, so its model norm is one and $g_\alpha=0$.
The local trace then vanishes over every tube avoiding $\mathcal Z$.
This concerns the full model support, including possible vertical
cohomology, not just the closure of its generic support.
\end{proof}

\paragraph*{A support qualification that must not be omitted.}
For an arbitrary complex, coefficients of $z_C$ can cancel even when
its cohomology support is larger.  The theorem produces a form on
$X\setminus Z$, and does not automatically extend it through those
canceled components.  To obtain the exact final support of a
\emph{prescribed cycle}, choose the complexes component by component
as in the next corollary. The componentwise construction avoids
attempting to extend a Green form across canceled components.

\subsection{Arbitrary prescribed cycles}

\begin{corollary}[Every prescribed cycle has a strong model-curvature pair]
\label{r11p:any-cycle-pair}
Let $X/K$ be smooth geometrically integral projective and let
$Z=\sum_j a_j[Z_j]$ be a rational codimension-$r$ cycle, written with
distinct integral components and nonzero coefficients.  After choosing
projective model data, there are an actual model Chern form $\Omega_Z$
and an actual
\[
                  G_Z\in\mathcal B^{r-1,r-1}
                           (X^{\rm an}\setminus|Z|^{\rm an})
\]
satisfying $D[G_Z]_\delta+\delta_Z=[\Omega_Z]_\delta$ on all compact
actual tests.  No generic cycle representative is moved.
\end{corollary}

\begin{proof}
On the regular projective generic fiber, $\mathcal O_{Z_j}$ has a
finite globally locally free resolution.  Its degree-$r$ supported
character is $[Z_j]$, including the leading coefficient; this is
Stacks 0FEW, or the leading-term part of localized Riemann--Roch.
It is not necessary for $Z_j$ to be a regular immersion.

These finitely many generic bundle complexes extend to a common
projective model after a vertical modification.  For completeness,
extend each bundle coherently, twist to obtain finitely many generators,
and close the graph of its generic quotient in a Grassmann bundle.
Its tautological quotient is a model vector bundle with the original
generic fiber.  Use the product of the finitely many Grassmann bundles.
Every generic differential extends after multiplying by a common
power of a uniformizer; the resulting successive products are zero
because they vanish on the dense generic fiber.  The scaled generic
complex is chain-isomorphic to the original one by degreewise scalar
rescaling.  Its full model cohomology support has generic support
exactly $Z_j$.

Apply Theorem~\ref{r11p:complex-green} to each complex and sum with
coefficients $a_j$.  The union of the removed sets is exactly
$|Z|$, because the prescribed cycle was first written with distinct
nonzero components.  The zero cycle is handled by the zero form.
\end{proof}

\subsection{Global primitives for generic Chow-zero model characters}

\begin{corollary}[Generically acyclic complexes]
\label{r11p:acyclic}
If a bounded model complex is exact on all of $X$, then
$\sum_i(-1)^i\operatorname{ch}_r(\overline C_i)$ has a global actual
$D$-primitive on $X^{\rm an}$.
\end{corollary}

\begin{proof}
The same relative construction is made with generic open $U=X$;
its final generic multisection is finite flat over all of $X$ and all
generic line divisors are empty.  Formula \eqref{r11p:complex-green-formula}
is global, and \eqref{r11p:complex-green-equation} has zero residue.
Actual current injectivity identifies its $D$ with the displayed
actual curvature.
\end{proof}

\begin{proposition}[A generic $K_0$ relation has an actual primitive]
\label{r11p:generic-k-primitive}
Let $\xi\in K_0(\mathcal X)_{\mathbb Q}$.  If $\xi|_X=0$ in
$K_0(X)_{\mathbb Q}$, then every positive-degree actual
$\operatorname{ch}_r(\xi)$ has a global actual $D$-primitive.
\end{proposition}

\begin{proof}
A relation in $K_0(X)$ is a finite combination of short exact sequence
relations and differences of isomorphic vector bundles.  Extend the
finitely many bundles to a common vertical model refinement, which
does not change any original model metric on the generic fiber.
For a generic exact sequence $0\to A\to B\to C\to0$, close the graph
of the subbundle $A\subset B$ in a Grassmann bundle of the already
chosen model extension of $B$.  This produces a model exact sequence
\[
                         0\to S\to B\to Q\to0
\]
with $S|_X\simeq A$ and $Q|_X\simeq C$.  Strict model Whitney gives
\[
      \operatorname{ch}_r(B-A-C)
       =\operatorname{ch}_r(S-A)+\operatorname{ch}_r(Q-C).
\]
After clearing uniformizer powers, the two generic isomorphisms are
model maps giving two-term generically acyclic complexes.  Their
global primitives exist by Corollary~\ref{r11p:acyclic}.
The same handles the isomorphic-bundle differences in the chosen
finite presentation of the $K_0$ relation.  Add the primitives and
divide by the cleared integer denominator.
\end{proof}

\begin{proposition}[Only the generic degree-$r$ Chow class matters]
\label{r11p:generic-chow-primitive}
If $\xi\in K_0(\mathcal X)_{\mathbb Q}$ satisfies
$\operatorname{ch}_r(\xi|_X)=0$ in
$\operatorname{CH}^r(X)_{\mathbb Q}$, then its actual model
$\operatorname{ch}_r(\xi)$ has a global actual $D$-primitive.
\end{proposition}

\begin{proof}
Let $\psi^2$ be the Adams operation on model vector-bundle $K_0$ and set
\[
 e_r(T)=\prod_{\substack{0\leq j\leq n\\j\ne r}}
                  \frac{T-2^j}{2^r-2^j},
 \qquad
                            \xi_r=e_r(\psi^2)\xi.
\]
On the smooth projective generic fiber the rational Chern character
$K_0(X)_{\mathbb Q}\simeq\bigoplus_{j=0}^n
\operatorname{CH}^j(X)_{\mathbb Q}$ is an isomorphism, and
$\operatorname{ch}_j(\psi^2\zeta)=2^j\operatorname{ch}_j(\zeta)$.
All generic Chern components of $\xi_r$ are zero, so $\xi_r|_X=0$.

The same Adams identity is strict for actual model Chern forms:
it is a universal symmetric-root identity, hence follows from
the strict operational identities and flag-section descent in
Section~\ref{h22c:section}.
Therefore $\operatorname{ch}_r(\xi_r)=\operatorname{ch}_r(\xi)$
as actual forms.  Apply Proposition~\ref{r11p:generic-k-primitive}
to $\xi_r$.
\end{proof}



\section{The supported-complex construction over a nongeometrically integral base}
\label{r12i:section}

\begin{lemma}[Native integral scope of the supported-complex construction]
\label{r12i:integral-PC}
Keep the original finite extension $K/\mathbf Q_p$ and its absolute
value.  The supported-complex theorem of Section~\ref{r11p:supported-complex}, the exact-support
Green-pair construction for every prescribed rational cycle, and
the generic-Chow-zero global primitive construction apply to every
smooth integral projective $K$-variety $X$.  Geometric integrality
over $K$ is not required.  Every actual form, test and finite-trace
operation in these statements is in the original $K$-analytic theory.
\end{lemma}

\begin{proof}
The supported-complex section and its main theorem already assume
only that $X$ is smooth and integral.  We check the one potentially
sensitive ingredient and then the passage to an arbitrary prescribed
cycle, whose preceding corollary was stated with the stronger hypothesis.

Let $F\to Y$ be the smooth projective flag fibration appearing after
the kernel-Grassmann modification and its characteristic-zero generic
resolution.  Both $F$ and $Y$ are smooth integral $K$-varieties.
Put $d=\dim F-\dim Y$.  If $d=0$, the map is already finite etale
and its integral source can be used.  If $d>0$, the uniform fiber
dimension argument in Lemma~\ref{r11p:multisection} gives a nonempty open of tuples of
high-degree very ample sections whose common zero locus has finite
nonempty fibers over every point of $Y$.  Its dimension estimate is
made on every geometric fiber and does not require the total space
or the base to be geometrically integral.

Intersect this parameter open with the smooth complete-intersection
Bertini open.  The latter requires smoothness of the source and a
linear system giving an immersion, not geometric irreducibility;
these are exactly the hypotheses of Stacks Lemma 33.47.3,
Tag~\texttt{0FD6}.  Successive restrictions of the very ample linear
system still give immersions.  Equivalently, after algebraic base
change the smoothness conditions are imposed simultaneously on the
finitely many geometric components; the resulting nonempty open is
defined over $K$.  Since $K$ is infinite it has a $K$-rational tuple.
No analytic scalar extension is involved in this algebraic check.

The resulting complete intersection $T_0^{\rm all}$ is smooth,
projective, and of pure dimension $\dim Y$.  Its map to $Y$ is
proper with finite nonempty fibers, hence finite surjective.
Miracle flatness applies, since its source is Cohen--Macaulay and
$Y$ is regular; see Stacks Tag~\texttt{00R4}.
Every connected component of a regular scheme is integral, because
its irreducible components are disjoint.  Choose any nonempty
connected component $T_0$.  It is open and closed in
$T_0^{\rm all}$, so its finite map is still flat.  Its image is both
open and closed in integral $Y$, and therefore all of $Y$.
Consequently $T_0\to Y$ is smooth-source, integral, finite flat and
surjective, with its own native positive degree.

Choosing this component does not spoil the supported $K_0$ identity:
the identity is pulled back to the chosen component, rather than
being an equality that would require summing the classes of all
components of a fixed intersection cycle.  Subsequent target-flag
graph closures and generic resolutions are birational on this
integral source and preserve its full prescribed finite open.
The leading supported-character identity then uses
$f_*[T]=d[X]$ with this actual function-field degree $d$.
Normality needed by the minimal trace follows from smoothness of
the generic source and target.  None of these statements requires
geometric integrality.

For a prescribed cycle $Z=\sum_j a_j[Z_j]$ on this $X$, resolve each
$\mathcal O_{Z_j}$ by finitely many generic vector bundles.  Such
resolutions use regularity and the projective resolution property,
which hold for smooth integral $X$ over $K$.  Extend the bundles by
the same finite Grassmann-graph construction to a projective flat
model, and extend the differentials after uniformizer rescaling.
Each model complex has generic support exactly $Z_j$.  The preceding
supported-complex theorem applies component by component.  Summing
with the nonzero coefficients gives a Green pair on exactly
$X^{\rm an}\setminus|Z|^{\rm an}$, with the original cycle and all
actual tests retained.  This proves the stated extension of the
exact-support corollary.

The global primitive proof uses the same construction with generic
open $U=X$.  Generic $K_0$ relations, Grassmann subbundle graphs,
and rational Adams projectors require smooth integral projectivity,
not geometric integrality.  Their primitives consequently retain
the original global scalar-curvature line sums and finite-flat maps
over $K$, as in Section~\ref{r11p:supported-complex}.
\end{proof}

\begin{remark}[Finite constant extensions are treated as native schemes]\label{h22:native-constant-extensions}
Let $K'/K$ be finite and let $Y/K'$ be smooth integral projective.
The composite $Y\to\operatorname{Spec}K$ is smooth because
$K'/K$ is separable, and it is still integral and projective.
A projective flat $R'$-model of $Y$ is projective and flat over $R$.
Lemma~\ref{r12i:integral-PC} can therefore be run directly on this
$R$-model, producing $K$-actual forms and testing against arbitrary
$K$-actual tests.  This is not an inference from a formula tested
only over $K'$.

The special fiber over $R$ keeps its ramification thickening, and
the degree of a finite generic map is its full degree as a
$K$-morphism.  Thus no comparison of native skeleton measures over
different ground fields is required for this scope clarification.
The present lemma makes this native treatment explicit for the
exact-support corollary.
\end{remark}

\begin{thebibliography}{99}
\bibitem{MihLT} A. Mihatsch, \emph{$\delta$-Forms on Lubin--Tate spaces},
arXiv:2112.10018v4, 7 September 2024.
\url{https://arxiv.org/abs/2112.10018v4}.
\bibitem{MihTrop} A. Mihatsch, \emph{On Tropical Intersection Theory},
arXiv:2107.12067v3, 28 August 2022.
\url{https://arxiv.org/abs/2107.12067v3}.
\bibitem{CLD} A. Chambert-Loir and A. Ducros,
\emph{Formes diff\'erentielles r\'eelles et courants sur les espaces de Berkovich},
arXiv:1204.6277v2. Section and proposition numbers in this paper refer
to the PDF version. \url{https://arxiv.org/pdf/1204.6277v2}.
\bibitem{GM} W. Gubler and F. Martin, \emph{On Zhang's semipositive metrics},
Documenta Mathematica 24 (2019), 331--372; arXiv:1608.08030v3.
\url{https://arxiv.org/abs/1608.08030v3}.
\bibitem{Fulton} W. Fulton, \emph{Intersection Theory}, second edition,
Springer, 1998, Corollary 18.1.2.
\bibitem{Stacks} The Stacks Project, Tags 0D4I, 0FD6, 0G4C, 00R4,
0FB3, 0FBE, 0FET, 0FEW, 0FEY and 0FEI; accessed 23 November 2026.
\url{https://stacks.math.columbia.edu/}.
\bibitem{R10R12} OpenMath Strong Green project,
\emph{Companion foundations, R10--R12}, frozen edition v0.5.1.
\url{https://openmath-atlas.org/projects/strong-green/editions/v0.5.1/}.
\end{thebibliography}
\end{document}
